\documentclass[12pt]{article}

\usepackage[margin=1in]{geometry}
\usepackage{amsmath,amssymb}
\usepackage[most]{tcolorbox}

\everymath{\displaystyle}

\sloppy
\allowdisplaybreaks
\widowpenalty=10000
\clubpenalty=10000

\newcounter{probnum}

\definecolor{LightGreen}{RGB}{144,238,144}


\newenvironment{problembox}{\par\noindent\rule{\linewidth}{0.5pt}\vspace{5pt}\par}{\vspace{2pt}}

\newtcolorbox{notebox}{
  enhanced,
  breakable,
  colback=LightGreen!20,
  colframe=LightGreen!60!black,
  boxrule=0.8pt,
  arc=0pt,
  left=8pt, right=8pt, top=6pt, bottom=6pt,
  before skip=12pt, after skip=29pt,
  boxsep=1pt
}

\setlength{\parindent}{0pt}
\setlength{\parskip}{0.5em}

\title{Collection of Real Analysis Problems}
\author{Zeta Zero}
\date{}

\begin{document}
\maketitle

\begin{notebox}
For corrections or to suggest a new problem for this collection, please email \texttt{category.theory.for.cats@gmail.com}.
\end{notebox}

\stepcounter{probnum}
\begin{problembox}
\textbf{\theprobnum.} Let $f_n : \mathbb{R} \to \mathbb{R}$ ($n \ge 1$) be the sequence of functions with general term
\[
f_n(x) = \frac{n^2 x}{x^2 + n^2}.
\]
a) Prove that the sequence $(f_n)$ converges pointwise on $\mathbb{R}$ and determine its pointwise limit $f$.\\
b) Show that $(f_n)$ does not converge uniformly on $\mathbb{R}$.\\
c) Show that $\displaystyle\int_0^x f(t)\,dt = \lim_{n \to \infty} \int_0^x f_n(t)\,dt$ for every $x \in \mathbb{R}$.
\end{problembox}

\stepcounter{probnum}
\begin{problembox}
\textbf{\theprobnum.} Given a nonempty set $A \subseteq \mathbb{R}$, a function $f : A \to \mathbb{R}$ is called a \emph{Baire function} if there exists a sequence of functions $f_n : A \to \mathbb{R}$ ($n \ge 1$), continuous on $A$, which converges pointwise to $f$ on $A$.\\
a) Prove that if $f : \mathbb{R} \to \mathbb{R}$ is differentiable on $\mathbb{R}$, then its derivative $f'$ is a Baire function.\\
b) Give an example of a Baire function that is not the derivative of any function.
\end{problembox}

\stepcounter{probnum}
\begin{problembox}
\textbf{\theprobnum.} Let $f_n : [a,b] \to \mathbb{R}$ ($n \ge 1$) be a sequence of functions satisfying:\\
(i) all functions $f_n$ are differentiable on $[a,b]$;\\
(ii) the sequence $(f_n)$ converges pointwise on $[a,b]$ to a differentiable function $f : [a,b] \to \mathbb{R}$;\\
(iii) the sequence $(f_n')$ converges pointwise on $[a,b]$ to a continuous function $g : [a,b] \to \mathbb{R}$.\\
Prove that $f'(x) = g(x)$ for all $x \in [a,b]$.
\end{problembox}

\stepcounter{probnum}
\begin{problembox}
\textbf{\theprobnum.} Let $g : [0,1] \to \mathbb{R}$ be a continuous function with $g(1) = 0$, and let $f_n : [0,1] \to \mathbb{R}$ ($n \ge 1$) be the sequence of functions defined by $f_n(x) := x^n g(x)$ for every $n \ge 1$ and every $x \in [0,1]$. Prove that the sequence $(f_n)$ converges uniformly on $[0,1]$.
\end{problembox}

\stepcounter{probnum}
\begin{problembox}
\textbf{\theprobnum.} Let $A \subseteq \mathbb{R}$ be a nonempty set, and let $f_n : A \to \mathbb{R}$ ($n \ge 1$) be a sequence of uniformly continuous functions which converges uniformly to $f : A \to \mathbb{R}$. Prove that $f$ is uniformly continuous.
\end{problembox}

\stepcounter{probnum}
\begin{problembox}
\textbf{\theprobnum.} Let $f_n : \mathbb{R} \to \mathbb{R}$ ($n \ge 1$) be a sequence of functions that converges uniformly on $\mathbb{R}$ to the function $f : \mathbb{R} \to \mathbb{R}$. Suppose that for each $n \ge 1$ there exists $\ell_n = \lim_{x \to \infty} f_n(x) \in \mathbb{R}$. Prove that the limits $\lim_{n \to \infty} \ell_n$ and $\lim_{x \to \infty} f(x)$ both exist and are equal.
\end{problembox}

\stepcounter{probnum}
\begin{problembox}
\textbf{\theprobnum.} Let $f_1 : [a,b] \to \mathbb{R}$ be a Riemann integrable function, and let $f_n : [a,b] \to \mathbb{R}$ ($n \ge 1$) be the sequence of functions defined recursively by
\[
f_{n+1}(x) = \int_a^x f_n(t)\,dt \quad \text{for every } x \in [a,b] \text{ and every } n \ge 1.
\]
Prove that the sequence $(f_n)$ converges uniformly to the zero function on $[a,b]$.
\end{problembox}

\stepcounter{probnum}
\begin{problembox}
\textbf{\theprobnum.} Let $f_n : [0,1] \to \mathbb{R}$ ($n \ge 1$) be the sequence defined recursively by $f_1(x) = 0$,
\[
f_{n+1}(x) = f_n(x) + \frac{1}{2}\left(x - f_n^2(x)\right) \quad \text{for every } x \in [0,1].
\]
Prove that the sequence $(f_n)$ converges uniformly on $[0,1]$ to the function $f(x) = \sqrt{x}$.
\end{problembox}

\stepcounter{probnum}
\begin{problembox}
\textbf{\theprobnum.} Calculate
\[
\lim_{n \to \infty} \int_0^1 \sqrt{\frac{x}{2} + \sqrt{\frac{x}{2} + \cdots + \sqrt{\frac{x}{2} + \sqrt{2}}}} \, dx
\]
with $n$ radicals in the nested expression.
\end{problembox}

\stepcounter{probnum}
\begin{problembox}
\textbf{\theprobnum.} Let $f : \mathbb{R} \to \mathbb{R}$ be a function. Prove that there exists a sequence of polynomial functions $f_n : \mathbb{R} \to \mathbb{R}$ ($n \ge 1$), uniformly convergent to $f$, if and only if $f$ is a polynomial function.
\end{problembox}

\stepcounter{probnum}
\begin{problembox}
\textbf{\theprobnum.} Let $m \in \mathbb{N}$ and let $f_n : [0,1] \to \mathbb{R}$ ($n \ge 1$) be a sequence of polynomial functions of degree at most $m$, which converge pointwise to the function $f : [0,1] \to \mathbb{R}$. Prove that $f$ is a polynomial function of degree at most $m$, and that the sequence $(f_n)$ converges uniformly to $f$ on $[0,1]$.
\end{problembox}

\stepcounter{probnum}
\begin{problembox}
\textbf{\theprobnum.} Determine the set of convergence of the series of functions
\[
\sum_{n \ge 1} \left(\frac{2n^2+5}{7n^2+3n+2}\right) \left(\frac{x}{2x+1}\right)^n, \qquad x \in \mathbb{R} \setminus \left\{-\tfrac{1}{2}\right\}.
\]
\end{problembox}

\stepcounter{probnum}
\begin{problembox}
\textbf{\theprobnum.} Define the sequence $x_1 = a \in \mathbb{R}$; $x_{n+1} = x_n^2 + x_n + 1$. Calculate
\[
\lim_{n \to \infty} \frac{(x_{n+1}+1)^2}{(x_1^2+2)(x_2^2+2)\cdots(x_n^2+2)}.
\]
\end{problembox}

\stepcounter{probnum}
\begin{problembox}
\textbf{\theprobnum.} Consider the sequence $(x_n)_{n \ge 0}$ defined by $x_0 > 0$; $x_{n+1} = \dfrac{x_n}{1+n x_n^2}$. Study the convergence of the sequences $(x_n)_{n\ge 0}$ and $(n x_n)_{n \ge 0}$, and determine their possible limits.
\end{problembox}

\stepcounter{probnum}
\begin{problembox}
\textbf{\theprobnum.} Consider the sequence $(x_n)$ of real numbers, defined by $x_1 > 0$; $x_1 + x_1^2 < 1$ and $x_{n+1} = x_n + \dfrac{x_n^2}{n^2}$ for every $n \ge 1$. Show that the sequences $\left(\dfrac{1}{x_n} - \dfrac{1}{n-1}\right)$ and $(x_n)$ are convergent.
\end{problembox}

\stepcounter{probnum}
\begin{problembox}
\textbf{\theprobnum.} Let $(a_n)_{n \ge 1}$ be a sequence with strictly positive terms. If
\[
a_n \le a_{2n} + a_{2n+1}, \qquad \forall n \in \mathbb{N}^*,
\]
study the convergence of the series $\displaystyle\sum_{n=1}^{\infty} a_n$.
\end{problembox}

\stepcounter{probnum}
\begin{problembox}
\textbf{\theprobnum.} Let $f : \mathbb{R} \to \mathbb{R}$ satisfy
\[
f(x) + f''(x) = -x\, g(x)\, f'(x),
\]
where $g : \mathbb{R} \to \mathbb{R}$ satisfies $g(x) > 0$ for all $x \in \mathbb{R}$. Show that $f$ is bounded.
\end{problembox}

\stepcounter{probnum}
\begin{problembox}
\textbf{\theprobnum.} Consider a sequence $(u_n)_{n \ge 1}$ defined by initial terms $u_1, u_2 \in (0,1)$ and the recursive relation
\[
u_{n+2} = \frac{1}{n} u_{n+1}^2 + \frac{n-1}{n} \sqrt{u_n} \quad \text{for all } n \in \mathbb{N}.
\]
Determine the limit of the sequence $(u_n)$ as $n \to \infty$.
\end{problembox}

\stepcounter{probnum}
\begin{problembox}
\textbf{\theprobnum.} Let the sequence $(u_n)$ be defined by $u_1 = \dfrac{1}{2}$ and
\[
u_{n+1} = u_n + \frac{u_n^2}{n^2}, \qquad \forall n \in \mathbb{N}^*.
\]
a) Prove that $u_n + n^2 \ge \dfrac{n(n+1)}{\sqrt{2}}$.\\
b) Prove that $(u_n)$ has a finite limit.
\end{problembox}

\stepcounter{probnum}
\begin{problembox}
\textbf{\theprobnum.} a) Calculate $\displaystyle\int_0^1 x \sin(\pi x^2)\,dx$.\\
b) Calculate $\displaystyle\lim_{n \to \infty} \frac{1}{n} \sum_{k=1}^{n-1} k \int_{k/n}^{(k+1)/n} \sin(\pi x^2)\,dx$.
\end{problembox}

\stepcounter{probnum}
\begin{problembox}
\textbf{\theprobnum.} Let $f : [0,1] \to (0,\infty)$ be a continuous function.\\
a) Prove that for every integer $n \ge 1$ there exists a unique partition $(a_0, a_1, \ldots, a_n)$ of $[0,1]$ such that
\[
\int_{a_k}^{a_{k+1}} f(t)\,dt = \frac{1}{n} \int_0^1 f(t)\,dt \quad \text{for all } k \in \{0,1,\ldots,n-1\}.
\]
b) For each integer $n \ge 1$ define $\bar{a}_n = \dfrac{a_1 + \cdots + a_n}{n}$, where $(a_0,a_1,\ldots,a_n)$ is the unique partition from part a). Show that the sequence $(\bar{a}_n)_{n \ge 1}$ is convergent and determine its limit.
\end{problembox}

\stepcounter{probnum}
\begin{problembox}
\textbf{\theprobnum.} Let $f : [0,1] \to (0,\infty)$ be a continuous function. For each $n \in \mathbb{N}$, $n \ge 2$, consider the partition $(t_0,t_1,\ldots,t_n)$ of $[0,1]$ with the property that
\[
\int_{t_0}^{t_1} f(t)\,dt = \int_{t_1}^{t_2} f(t)\,dt = \cdots = \int_{t_{n-1}}^{t_n} f(t)\,dt.
\]
Determine $\displaystyle\lim_{n \to \infty} n \left(\frac{1}{f(t_1)} + \frac{1}{f(t_2)} + \cdots + \frac{1}{f(t_n)}\right)$.
\end{problembox}

\stepcounter{probnum}
\begin{problembox}
\textbf{\theprobnum.} Let $f : [a,b] \to \mathbb{R}$ be a differentiable function with Riemann integrable derivative. For each $n \in \mathbb{N}$ denote
\[
A_n := \frac{b-a}{n} \sum_{k=0}^{n-1} f\left(a + k\,\frac{b-a}{n}\right), \qquad B_n := \frac{b-a}{n} \sum_{k=1}^{n} f\left(a + k\,\frac{b-a}{n}\right).
\]
Prove that
\[
\lim_{n \to \infty} n \left(A_n - \int_a^b f(x)\,dx\right) = \frac{b-a}{2}\big(f(a) - f(b)\big)
\]
and
\[
\lim_{n \to \infty} n \left(B_n - \int_a^b f(x)\,dx\right) = \frac{b-a}{2}\big(f(b) - f(a)\big).
\]
\end{problembox}

\stepcounter{probnum}
\begin{problembox}
\textbf{\theprobnum.} Determine $\displaystyle\lim_{n \to \infty} n \left(\frac{1}{n+1} + \frac{1}{n+2} + \cdots + \frac{1}{2n} - \ln 2\right)$.
\end{problembox}

\stepcounter{probnum}
\begin{problembox}
\textbf{\theprobnum.} Let $f : [0,1] \to \mathbb{R}$ be an increasing, twice differentiable function, with continuous second derivative on $[0,1]$, and let $(A_n)$, $(B_n)$ be the sequences defined by
\[
A_n := \frac{1}{n} \sum_{k=0}^{n-1} f\left(\frac{k}{n}\right), \qquad B_n := \frac{1}{n} \sum_{k=1}^{n} f\left(\frac{k}{n}\right).
\]
The interval $[A_n, B_n]$ is divided into three equal segments. Prove that for $n$ sufficiently large, the number $I := \displaystyle\int_a^b f(x)\,dx$ belongs to the middle segment.
\end{problembox}

\stepcounter{probnum}
\begin{problembox}
\textbf{\theprobnum.} Let $f : [0,1] \to \mathbb{R}$ be a differentiable function with continuous derivative, and let
\[
s_n = \sum_{k=1}^{n} f\left(\frac{k}{n}\right) \quad \text{for every } n \in \mathbb{N}.
\]
Prove that the sequence $(s_{n+1} - s_n)_{n \ge 1}$ converges to $\displaystyle\int_0^1 f(x)\,dx$.
\end{problembox}

\stepcounter{probnum}
\begin{problembox}
\textbf{\theprobnum.} Let $f,g : [a,b] \to \mathbb{R}$ be functions with the property that there exist $\alpha \ge 0$ and $p \ge 1$ such that
\[
|f(x) - f(x')| \le \alpha\, |g(x) - g(x')|^p \quad \text{for all } x, x' \in [a,b].
\]
Prove that if $g$ is Riemann integrable, then $f$ is also Riemann integrable.
\end{problembox}

\stepcounter{probnum}
\begin{problembox}
\textbf{\theprobnum.} Let $f : [0,1] \to \mathbb{R}$ be a Riemann integrable function and let $p \ge 1$. Prove that the limit
\[
\lim_{n \to \infty} \frac{1}{n} \sum_{j=1}^{n} \left| f\left(\frac{j-1}{n}\right) - f\left(\frac{j}{n}\right) \right|^p
\]
exists, and determine its value.
\end{problembox}

\stepcounter{probnum}
\begin{problembox}
\textbf{\theprobnum.} (Riemann's function) Prove that the function $f : [0,1] \to \mathbb{R}$, defined by
\[
f(x) := \begin{cases} 1 & \text{if } x = 0, \\ 1/q & \text{if } x \in [0,1] \cap \mathbb{Q},\ x = p/q,\ p,q \in \mathbb{N},\ \gcd(p,q)=1, \\ 0 & \text{if } x \in [0,1] \setminus \mathbb{Q}, \end{cases}
\]
is Riemann integrable. (This function is also known in the literature as Thomae's function, the popcorn function, or the raindrop function.)
\end{problembox}

\stepcounter{probnum}
\begin{problembox}
\textbf{\theprobnum.} Prove that if the function $f : [a,b] \to \mathbb{R}$ has a finite limit at every point of $[a,b]$, then it is bounded.
\end{problembox}

\stepcounter{probnum}
\begin{problembox}
\textbf{\theprobnum.} Prove that the set of first-kind discontinuity points of a function $f : I \to \mathbb{R}$, where $I \subseteq \mathbb{R}$ is an interval, is at most countable.
\end{problembox}

\stepcounter{probnum}
\begin{problembox}
\textbf{\theprobnum.} Prove that if the function $f : [a,b] \to \mathbb{R}$ has a finite limit at every point of $[a,b]$, then it is Riemann integrable.
\end{problembox}

\stepcounter{probnum}
\begin{problembox}
\textbf{\theprobnum.} Prove that if $f : [0,1] \to \mathbb{R}$ is a Riemann integrable function, then the function $g : [0,1] \to \mathbb{R}$, defined by $g(x) = f(x^2)$, is also Riemann integrable.
\end{problembox}

\stepcounter{probnum}
\begin{problembox}
\textbf{\theprobnum.} Let $f : [0,1] \to \mathbb{R}$ be a Riemann integrable function with the property that $0 < \left|\displaystyle\int_0^1 f(x)\,dx\right| \le 1$. Prove that there exist $x_1, x_2 \in [0,1]$ such that $x_1 \ne x_2$ and
\[
\int_{x_1}^{x_2} f(x)\,dx = (x_1 - x_2)^{2002}.
\]
\end{problembox}

\stepcounter{probnum}
\begin{problembox}
\textbf{\theprobnum.} Let $f : [0,1] \to \mathbb{R}$ be an increasing function and let $n \ge 1$ be an integer. Prove that
\[
\int_0^1 f(x)\,dx \le (n+1) \int_0^1 x^n f(x)\,dx.
\]
Specify all continuous increasing functions $f$ for which equality holds.
\end{problembox}

\stepcounter{probnum}
\begin{problembox}
\textbf{\theprobnum.} Let $f : \left[0, \dfrac{\pi}{2}\right] \to \mathbb{R}$ be a continuous decreasing function. Prove that
\[
\int_{\pi/2 - 1}^{\pi/2} f(x)\,dx \le \int_0^{\pi/2} f(x) \cos x\,dx \le \int_0^1 f(x)\,dx.
\]
For which functions $f$ do both inequalities hold with equality?
\end{problembox}

\stepcounter{probnum}
\begin{problembox}
\textbf{\theprobnum.} Let $f : [0,1] \to \mathbb{R}$ be a differentiable function and let $M \ge 0$ be such that $|f'(x)| \le M$ for all $x \in [0,1]$. Prove that
\[
\int_0^1 f^2(x)\,dx - \left(\int_0^1 f(x)\,dx\right)^2 \le \frac{M^2}{12}.
\]
\end{problembox}

\stepcounter{probnum}
\begin{problembox}
\textbf{\theprobnum.} Let $f_1, \ldots, f_n : [0,1] \to (0,\infty)$ be continuous functions and let $\sigma$ be a permutation of $\{1,\ldots,n\}$. Prove that
\[
\prod_{i=1}^{n} \int_0^1 \frac{f_i^2(x)}{f_{\sigma(i)}(x)}\,dx \ge \prod_{i=1}^{n} \int_0^1 f_i(x)\,dx.
\]
\end{problembox}

\stepcounter{probnum}
\begin{problembox}
\textbf{\theprobnum.} Prove that for any continuous function $f : [-1,1] \to \mathbb{R}$ the inequality
\[
\int_{-1}^{1} f^2(x)\,dx \ge \frac{1}{2}\left(\int_{-1}^{1} f(x)\,dx\right)^2 + \frac{3}{2}\left(\int_{-1}^{1} x f(x)\,dx\right)^2
\]
holds. For which functions $f$ does equality hold?
\end{problembox}

\stepcounter{probnum}
\begin{problembox}
\textbf{\theprobnum.} a) Given continuous functions $u, v : [0,1] \to \mathbb{R}$, prove the Cauchy--Bunyakovsky--Schwarz inequality
\[
\left(\int_0^1 u(x) v(x)\,dx\right)^2 \le \int_0^1 u^2(x)\,dx \int_0^1 v^2(x)\,dx.
\]
b) Let $\mathcal{C}$ be the set of all differentiable functions $f : [0,1] \to \mathbb{R}$, with continuous derivative $f'$ on $[0,1]$ and with $f(0)=0$, $f(1)=1$. Determine
\[
\min_{f \in \mathcal{C}} \int_0^1 \sqrt{1+x^2}\,(f'(x))^2\,dx,
\]
as well as all functions $f \in \mathcal{C}$ for which this minimum is attained.
\end{problembox}

\stepcounter{probnum}
\begin{problembox}
\textbf{\theprobnum.} Let $f : [0,\infty) \to \mathbb{R}$ be a continuous function such that
\[
\int_0^n f(x) f(n-x)\,dx = \int_0^n f^2(x)\,dx
\]
for every natural number $n \ge 1$. Prove that $f$ is periodic.
\end{problembox}

\stepcounter{probnum}
\begin{problembox}
\textbf{\theprobnum.} Let $f : [0,1] \to \mathbb{R}$ be an integrable function such that
\[
\int_0^1 f(x)\,dx = \int_0^1 x f(x)\,dx = 1.
\]
Prove that $\displaystyle\int_0^1 f^2(x)\,dx \ge 4$.
\end{problembox}

\stepcounter{probnum}
\begin{problembox}
\textbf{\theprobnum.} Let $\mathcal{M}$ be the set of all continuous functions $f : [0,\pi] \to \mathbb{R}$, with the property that
\[
\int_0^{\pi} f(x) \sin x\,dx = \int_0^{\pi} f(x) \cos x\,dx = 1.
\]
Determine $\displaystyle\min_{f \in \mathcal{M}} \int_0^{\pi} f^2(x)\,dx$.
\end{problembox}

\stepcounter{probnum}
\begin{problembox}
\textbf{\theprobnum.} Let $p > 1$ be a real number and let $f : [a,b] \to [0,\infty)$ be a continuous function. Prove that
\[
\left(\int_a^b f(x)\,dx\right)^p \le (b-a)^{p-1} \int_a^b f(x)^p\,dx.
\]
\end{problembox}

\stepcounter{probnum}
\begin{problembox}
\textbf{\theprobnum.} Let $f : [a,b] \to \mathbb{R}$ be a Riemann integrable function. Prove that the function $\varphi : (0,\infty) \to \mathbb{R}$, defined by
\[
\varphi(r) := \left(\frac{1}{b-a} \int_a^b |f(x)|^r\,dx\right)^{1/r},
\]
is increasing.
\end{problembox}

\stepcounter{probnum}
\begin{problembox}
\textbf{\theprobnum.} Let $p > 1$ be a real number. For every continuous function $f : [0,1] \to \mathbb{R}$ we denote
\[
I(f) := \int_0^1 x^p |f(x)|\,dx - \int_0^1 x |f(x)|^p\,dx.
\]
Determine the maximum value that $I(f)$ can take.
\end{problembox}

\stepcounter{probnum}
\begin{problembox}
\textbf{\theprobnum.} Let $f : [0,1] \to \mathbb{R}$ and $g : [0,1] \to (0,\infty)$ be continuous functions, with $f$ increasing. Prove that for every $t \in [0,1]$ the inequality
\[
\int_0^t f(x) g(x)\,dx \int_0^1 g(x)\,dx \le \int_0^t g(x)\,dx \int_0^1 f(x) g(x)\,dx
\]
holds.
\end{problembox}

\stepcounter{probnum}
\begin{problembox}
\textbf{\theprobnum.} Let $f : [0,1] \to [0,\infty)$ be a decreasing function. Prove that
\[
\int_0^1 f(x)\,dx \int_0^1 x f^2(x)\,dx \le \int_0^1 f^2(x)\,dx \int_0^1 x f(x)\,dx.
\]
\end{problembox}

\stepcounter{probnum}
\begin{problembox}
\textbf{\theprobnum.} Prove that if $f : [0,1] \to [0,\infty)$ is Riemann integrable, then
\[
\int_0^1 f(x)\,dx \int_0^1 f^2(x)\,dx \le \int_0^1 f^3(x)\,dx.
\]
\end{problembox}

\stepcounter{probnum}
\begin{problembox}
\textbf{\theprobnum.} Let $f, g : [a,b] \to \mathbb{R}$ be Riemann integrable functions such that $m \le f(x) \le M$ for every $x \in [a,b]$ and $\displaystyle\int_a^b g(x)\,dx = 0$. Prove that
\[
\left|\int_a^b f(x) g(x)\,dx\right| \le \frac{M-m}{2} \int_a^b |g(x)|\,dx.
\]
\end{problembox}

\stepcounter{probnum}
\begin{problembox}
\textbf{\theprobnum.} Let $f : [0,1] \to \mathbb{R}$ be a differentiable function with continuous derivative on $[0,1]$. Prove that if $f(0)=0$ and $0 < f'(x) \le 1$ for all $x \in [0,1]$, then
\[
\left(\int_0^1 f(x)\,dx\right)^2 \ge \int_0^1 f^3(x)\,dx.
\]
Give an example of a function $f$ for which equality holds.
\end{problembox}

\stepcounter{probnum}
\begin{problembox}
\textbf{\theprobnum.} a) Let $f : [0,1] \to \mathbb{R}$ be a differentiable function with continuous derivative on $[0,1]$. Prove that if there exists $a \in [0,1]$ such that $f(a) \le 0$, then
\[
\int_0^1 f(t)\,dt \le \int_0^1 |f'(t)|\,dt.
\]
b) Let $g : [0,1] \to \mathbb{R}$ be a differentiable function with continuous derivative on $[0,1]$. Prove that for every $x \in [0,1]$ the inequality
\[
\int_0^1 \left(g(t) - |g'(t)|\right) dt \le g(x) \le \int_0^1 \left(g(t) + |g'(t)|\right) dt
\]
holds.
\end{problembox}

\stepcounter{probnum}
\begin{problembox}
\textbf{\theprobnum.} Let $f : [0,1] \to \mathbb{R}$ be a differentiable function with continuous derivative on $[0,1]$. Prove that if $\displaystyle\int_0^1 f(x)\,dx = 0$, then for every $\alpha \in (0,1)$ the inequality
\[
\left|\int_0^{\alpha} f(x)\,dx\right| \le \frac{1}{8} \max_{x \in [0,1]} |f'(x)|
\]
holds.
\end{problembox}

\stepcounter{probnum}
\begin{problembox}
\textbf{\theprobnum.} Let $f : [-1,1] \to \mathbb{R}$ be a differentiable function, with the property that there exists $a \in (0,1)$ such that $\displaystyle\int_{-a}^{a} f(x)\,dx = 0$. Denoting $M := \sup_{x \in [-1,1]} |f'(x)|$, prove that:\\
a) $\left|\displaystyle\int_{-1}^{1} f(x)\,dx\right| \le M(1-a)$.\\
b) $\displaystyle\int_{-1}^{1} |f(x)|\,dx \ge 2|f(0)| - M$.
\end{problembox}

\stepcounter{probnum}
\begin{problembox}
\textbf{\theprobnum.} Let $f : [0,1] \to \mathbb{R}$ be a differentiable function such that $f(0)=f(1)=0$ and $|f'(x)| \le 1$ for all $x \in [0,1]$. Prove that
\[
\left|\int_0^1 f(x)\,dx\right| \le \frac{1}{4}.
\]
\end{problembox}

\stepcounter{probnum}
\begin{problembox}
\textbf{\theprobnum.} Let $f : [a,b] \to \mathbb{R}$ be a continuous function on $[a,b]$, twice differentiable with bounded second derivative on $(a,b)$, such that $f(a)=f(b)=0$. Prove that
\[
\int_a^b |f(x)|\,dx \le M \frac{(b-a)^3}{12}, \qquad \text{where } M := \sup\{|f''(x)| : x \in (a,b)\}.
\]
\end{problembox}

\stepcounter{probnum}
\begin{problembox}
\textbf{\theprobnum.} Let $f : [-1,1] \to \mathbb{R}$ be a three-times differentiable function with continuous third derivative on $[-1,1]$, such that $f(-1)=f(0)=f(1)=0$. Prove that
\[
\int_{-1}^{1} |f(x)|\,dx \le \frac{1}{12} \max_{x \in [-1,1]} |f'''(x)|.
\]
\end{problembox}

\stepcounter{probnum}
\begin{problembox}
\textbf{\theprobnum.} (A. Beurling's inequality) Let $f : [a,b] \to \mathbb{R}$ be twice differentiable with continuous second derivative on $[a,b]$. Prove that if $f(a)=f(b)=0$, then
\[
\int_a^b |f''(x)|\,dx \ge \frac{4}{b-a} \max_{x \in [a,b]} |f(x)|.
\]
\end{problembox}

\stepcounter{probnum}
\begin{problembox}
\textbf{\theprobnum.} Let $f : [0,1] \to \mathbb{R}$ be a differentiable function with continuous derivative on $[0,1]$. Prove that if $f(0)=f(1)=-\dfrac{1}{6}$, then
\[
\int_0^1 (f'(x))^2\,dx \ge 2 \int_0^1 f(x)\,dx + \frac{1}{4}.
\]
\end{problembox}

\stepcounter{probnum}
\begin{problembox}
\textbf{\theprobnum.} Let $f : [0,1] \to \mathbb{R}$ be a differentiable function with continuous derivative on $[0,1]$. Determine $f$, knowing that $f(1) = -\dfrac{1}{6}$ and
\[
\int_0^1 (f'(x))^2\,dx \le 2 \int_0^1 f(x)\,dx.
\]
\end{problembox}

\stepcounter{probnum}
\begin{problembox}
\textbf{\theprobnum.} Let $f : [0,1] \to \mathbb{R}$ be a differentiable function with continuous derivative on $[0,1]$. Prove that if $f\left(\dfrac{1}{2}\right) = 0$, then
\[
\int_0^1 (f'(x))^2\,dx \ge 12 \left(\int_0^1 f(x)\,dx\right)^2.
\]
\end{problembox}

\stepcounter{probnum}
\begin{problembox}
\textbf{\theprobnum.} Let $f : [a,b] \to \mathbb{R}$ be a differentiable function with continuous derivative on $[a,b]$. Prove that if $f(a)=0$, then
\[
\int_a^b f(x)^2\,dx \le (b-a)^2 \int_a^b f'(x)^2\,dx.
\]
\end{problembox}

\stepcounter{probnum}
\begin{problembox}
\textbf{\theprobnum.} Let $f : [-1,1] \to \mathbb{R}$ be a twice differentiable function with continuous second derivative on $[-1,1]$. Prove that if $f(0)=0$, then
\[
\int_{-1}^{1} (f''(x))^2\,dx \ge 10 \left(\int_{-1}^{1} f(x)\,dx\right)^2.
\]
\end{problembox}

\stepcounter{probnum}
\begin{problembox}
\textbf{\theprobnum.} Let $n \ge 0$ be an integer and let $f : [-1,1] \to \mathbb{R}$ be $2n+2$ times differentiable with continuous derivative of order $2n+2$ on $[-1,1]$. Prove that if $f(0)=f''(0)=\cdots=f^{(2n)}(0)=0$, then
\[
\int_{-1}^{1} \left(f^{(2n+2)}(x)\right)^2\,dx \ge (2n+2)!^2 (4n+5)^2 \left(\int_{-1}^{1} f(x)\,dx\right)^2.
\]
\end{problembox}

\stepcounter{probnum}
\begin{problembox}
\textbf{\theprobnum.} Prove that if $f : [0,1] \to \mathbb{R}$ is a convex function, then for every $n \in \mathbb{N}$ the inequality
\[
\int_0^1 f(x)\,dx \le \frac{1}{n+1}\left(f(0) + f\left(\frac{1}{n}\right) + f\left(\frac{2}{n}\right) + \cdots + f(1)\right)
\]
holds.
\end{problembox}

\stepcounter{probnum}
\begin{problembox}
\textbf{\theprobnum.} Let $f : [-1,1] \to \mathbb{R}$ be a continuous function, twice differentiable on $(-1,1)$, with the property that
\[
2 f'(x) + x f''(x) \ge 1 \quad \text{for all } x \in (-1,1).
\]
Prove that $\displaystyle\int_{-1}^{1} x f(x)\,dx \ge \frac{1}{3}$.
\end{problembox}

\stepcounter{probnum}
\begin{problembox}
\textbf{\theprobnum.} Show that $(\cos t)^p \le \cos(pt)$, for every $t \in \left[0, \dfrac{\pi}{2}\right]$ and $p \in (0,1)$.
\end{problembox}

\stepcounter{probnum}
\begin{problembox}
\textbf{\theprobnum.} Let $f : [0,1] \to \mathbb{R}$ be a function of class $C^1$, with $f(0)=0$. Show that
\[
\sup_{x \in [0,1]} |f(x)| \le \sqrt{\int_0^1 [f'(x)]^2\,dx}.
\]
\end{problembox}

\stepcounter{probnum}
\begin{problembox}
\textbf{\theprobnum.} Let $f : [1,+\infty) \to \mathbb{R}$ with $f(1)=1$ and
\[
f'(x) = \frac{1}{x^2 + [f(x)]^2}.
\]
Show that $\displaystyle\lim_{x \to \infty} f(x)$ exists and is smaller than $1 + \dfrac{\pi}{4}$.
\end{problembox}

\stepcounter{probnum}
\begin{problembox}
\textbf{\theprobnum.} Let $f$ and $g : [0,1] \to (0,+\infty)$ be two continuous functions satisfying the relation
\[
\sup_{x \in [0,1]} f(x) = \sup_{x \in [0,1]} g(x).
\]
Show that there exists $t \in [0,1]$ with $[f(t)]^2 + 3f(t) = [g(t)]^2 + 3g(t)$.
\end{problembox}

\stepcounter{probnum}
\begin{problembox}
\textbf{\theprobnum.} Let $y(h) = 1 - 2\sin^2(2\pi h)$ and $f(y) = \dfrac{2}{1+\sqrt{1-y}}$. Show that
\[
f(y(h)) = 2 - 4\sqrt{2\pi}\,|h| + O(h^2),
\]
where $\displaystyle\limsup_{h \to 0} \frac{O(h^2)}{h^2} < \infty$.
\end{problembox}

\stepcounter{probnum}
\begin{problembox}
\textbf{\theprobnum.} Let $f : [0,1] \to \mathbb{R}^+$ be a continuous function. Suppose there exist positive constants $m$ and $M$ such that $(f(x)-m)(f(x)-M) \le 0$ for all $x \in [0,1]$. Prove that
\[
\int_0^1 f(x)\,dx \int_0^1 \frac{1}{f(x)}\,dx \le \frac{(M+m)^2}{4Mm}.
\]
\end{problembox}

\stepcounter{probnum}
\begin{problembox}
\textbf{\theprobnum.} Let $f$ be a non-negative Riemann integrable function on the interval $[0,1]$ satisfying $\displaystyle\int_0^1 f(x)\,dx = 1$. Prove the inequality
\[
\int_0^1 \left(x - \int_0^1 u f(u)\,du\right)^2 f(x)\,dx \le \frac{1}{4}.
\]
\end{problembox}

\stepcounter{probnum}
\begin{problembox}
\textbf{\theprobnum.} Let $f : [1,2] \to \mathbb{R}$ be a differentiable function such that its derivative satisfies $|f'(x)| \le 1$ for all $x \in [1,2]$. Prove the inequality
\[
\left|\int_1^2 f(x)\,dx - \frac{f(1)+f(2)}{2}\right| \le \frac{1 - [f(1)-f(2)]^2}{4}.
\]
\end{problembox}

\stepcounter{probnum}
\begin{problembox}
\textbf{\theprobnum.} Let $f : \mathbb{R} \to \mathbb{R}$ be a function satisfying the Lipschitz condition $|f(x)-f(y)| \le M|x-y|$ for all $x,y \in \mathbb{R}$, where $M > 0$ is a constant. Prove that for any real numbers $a < b$:\\
(i) $\left|\dfrac{1}{b-a}\displaystyle\int_a^b f(x)\,dx - f\left(\dfrac{a+b}{2}\right)\right| \le M \dfrac{b-a}{4}$.\\
(ii) $\left|\dfrac{1}{b-a}\displaystyle\int_a^b f(x)\,dx - \dfrac{f(a)+f(b)}{2}\right| \le M \dfrac{b-a}{3}$.
\end{problembox}

\stepcounter{probnum}
\begin{problembox}
\textbf{\theprobnum.} Let $f : \mathbb{R} \to \mathbb{R}$ be a continuously differentiable function. Show that
\[
\left|f(1) - \int_0^1 f(x)\,dx\right| \le \frac{1}{2} \max_{x \in [0,1]} |f'(x)|.
\]
\end{problembox}

\stepcounter{probnum}
\begin{problembox}
\textbf{\theprobnum.} Find the smallest constant $C$ such that for every real polynomial $P(x)$ of degree $3$ that has a root in the interval $[0,1]$,
\[
\int_0^1 |P(x)|\,dx \le C \max_{x \in [0,1]} |P(x)|.
\]
\end{problembox}

\stepcounter{probnum}
\begin{problembox}
\textbf{\theprobnum.} Let $f : [0,1] \to [0,\infty)$ be a non-identically-zero, differentiable function with continuous derivative.\\
a) Prove that the sequence $a_n = \displaystyle\int_0^1 \frac{f(x)}{x^n+1}\,dx$ ($n \in \mathbb{N}^*$) is convergent.\\
b) If $a = \lim_{n \to \infty} a_n$, determine $\displaystyle\lim_{n \to \infty} \frac{a_n}{a^n}$.
\end{problembox}

\stepcounter{probnum}
\begin{problembox}
\textbf{\theprobnum.} Let $f : [0,1] \to [0,1]$ be an increasing function and let
\[
a_n = \int_0^1 \frac{1+f^n(x)}{1+f^{n+1}(x)}\,dx \qquad (n \ge 1).
\]
Prove that the sequence $(a_n)_{n \ge 1}$ is convergent and determine its limit.
\end{problembox}

\stepcounter{probnum}
\begin{problembox}
\textbf{\theprobnum.} a) Given an integer $n \ge 0$, calculate $\displaystyle\int_0^1 (1-t)^n e^t\,dt$.\\
b) Let $k \ge 0$ be an integer and let $(x_n)_{n \ge k}$ be the sequence with general term
\[
x_n = \sum_{i=k}^{n} \binom{i}{k} \left(e - 1 - \frac{1}{1!} - \frac{1}{2!} - \cdots - \frac{1}{i!}\right).
\]
Prove that $(x_n)_{n \ge k}$ is convergent and determine its limit.
\end{problembox}

\stepcounter{probnum}
\begin{problembox}
\textbf{\theprobnum.} Let $f : [0,1] \to \mathbb{R}$ be a continuous function and let $(a_n), (b_n)$ be sequences of real numbers with the property that
\[
\lim_{n \to \infty} \int_0^1 |f(x) - a_n x - b_n|\,dx = 0.
\]
Prove that:\\
a) The sequences $(a_n)$ and $(b_n)$ are convergent.\\
b) There exist $a, b \in \mathbb{R}$ such that $f(x) = ax+b$ for all $x \in \mathbb{R}$.
\end{problembox}

\stepcounter{probnum}
\begin{problembox}
\textbf{\theprobnum.} Let $f : [0,1] \to \mathbb{R}$ be a continuous function with the property that
\[
\int_0^1 f(x) g(x)\,dx = \int_0^1 f(x)\,dx \int_0^1 g(x)\,dx
\]
for every continuous, non-differentiable function $g : [0,1] \to \mathbb{R}$. Prove that $f$ is constant.
\end{problembox}

\stepcounter{probnum}
\begin{problembox}
\textbf{\theprobnum.} Calculate $\displaystyle\lim_{n \to \infty} \int_0^1 e^{x^n}\,dx$.
\end{problembox}

\stepcounter{probnum}
\begin{problembox}
\textbf{\theprobnum.} a) Prove that $\ln(1+x) \le x$ for all $x \ge 0$.\\
b) If $a > 0$, prove that $\displaystyle\lim_{n \to \infty} n \int_0^1 \frac{x^n}{a+x^n}\,dx = \ln\frac{a+1}{a}$.
\end{problembox}

\stepcounter{probnum}
\begin{problembox}
\textbf{\theprobnum.} Let $f : [0,1] \to \mathbb{R}$ be the function $f(x) = (x^2+1)e^x$. Calculate
\[
\lim_{n \to \infty} n \int_0^1 \left(f\left(\frac{x}{2^n}\right) - 1\right) dx.
\]
\end{problembox}

\stepcounter{probnum}
\begin{problembox}
\textbf{\theprobnum.} For each $n \in \mathbb{N}^*$ consider the function $f_n : [0,n] \to \mathbb{R}$, defined by $f_n(x) = \arctan([x])$, where $[x]$ denotes the integer part of $x$. Show that $f_n$ is Riemann integrable and determine
\[
\lim_{n \to \infty} \frac{1}{n} \int_0^n f_n(x)\,dx.
\]
\end{problembox}

\stepcounter{probnum}
\begin{problembox}
\textbf{\theprobnum.} Consider the continuous functions $f : [0,\infty) \to \mathbb{R}$ and $g : [0,1] \to \mathbb{R}$. If the limit $\lim_{x \to \infty} f(x) = L \in \mathbb{R}$ exists, prove that
\[
\lim_{n \to \infty} \frac{1}{n} \int_0^n f(x)\, g\left(\frac{x}{n}\right) dx = L \int_0^1 g(x)\,dx.
\]
\end{problembox}

\stepcounter{probnum}
\begin{problembox}
\textbf{\theprobnum.} For each $\alpha \in (0,1]$ denote
\[
I_n(\alpha) = \int_0^{\alpha} \ln\left(1+x+x^2+\cdots+x^{n-1}\right) dx, \qquad n \ge 2.
\]
Calculate:\\
a) $\displaystyle\lim_{n \to \infty} I_n(\alpha)$ for $\alpha \in (0,1)$;\\
b) $\displaystyle\lim_{n \to \infty} I_n(1)$.
\end{problembox}

\stepcounter{probnum}
\begin{problembox}
\textbf{\theprobnum.} Let $\alpha > 0$ and let $a, b \in \mathbb{R}$ be such that $a < b$. Determine
\[
\lim_{n \to \infty} \int_a^b \sqrt[n]{(x-a)^n + (b-x)^n + \alpha\big((x-a)(b-x)\big)^{n/2}}\,dx.
\]
\end{problembox}

\stepcounter{probnum}
\begin{problembox}
\textbf{\theprobnum.} Determine $\displaystyle\lim_{n \to \infty} n \int_{-1}^{0} (x+e^x)^n\,dx$.
\end{problembox}

\stepcounter{probnum}
\begin{problembox}
\textbf{\theprobnum.} Determine $\displaystyle\lim_{n \to \infty} n \int_0^1 (\cos x - \sin x)^n\,dx$.
\end{problembox}

\stepcounter{probnum}
\begin{problembox}
\textbf{\theprobnum.} Let $\{\epsilon_n\}_{n=1}^{\infty}$ be a sequence of positive real numbers such that $\lim_{n \to \infty} \epsilon_n = 0$. Evaluate the limit
\[
\lim_{n \to \infty} \frac{1}{n} \sum_{k=1}^{n} \ln\left(\frac{k}{n} + \epsilon_n\right).
\]
\end{problembox}

\stepcounter{probnum}
\begin{problembox}
\textbf{\theprobnum.} Let $\alpha, \beta$ be reals such that $2\beta \ge \alpha$ and suppose that they are not simultaneously negative. Let
\[
F(\alpha,\beta) = \lim_{n \to \infty} n^{\alpha} \int_0^{n^{-\beta}} x^{x+1}\,dx.
\]
Prove that $F(\alpha,\beta)$ exists and is finite. Hence find the value of $F(\alpha,\beta)$.
\end{problembox}

\stepcounter{probnum}
\begin{problembox}
\textbf{\theprobnum.} Let $\psi : \mathbb{R} \to \mathbb{R}$ be such that
\[
\int_x^{x+y} \psi(t)\,dt = \int_{x-y}^{x} \psi(t)\,dt.
\]
Find all possible $\psi(x)$.
\end{problembox}

\stepcounter{probnum}
\begin{problembox}
\textbf{\theprobnum.} (Wirtinger's inequalities)
\\
a) Let $f : [0,1] \to \mathbb{R}$ be a differentiable function with continuous derivative on $[0,1]$, with the property that $\displaystyle\int_0^1 f(x)\,dx = 0$. Prove that
\[
\int_0^1 f(x)^2\,dx \le \frac{1}{2} \int_0^1 f'(x)^2\,dx.
\]
\\
b) Let $f : [0,1] \to \mathbb{R}$ be differentiable and satisfy $f(0)=0$. Prove that
\[
\int_0^1 [f'(x)]^2\,dx \ge 3 \left(\int_0^1 f(x)\,dx\right)^2.
\]
\\
c) Let $f : [0,1] \to \mathbb{R}$ be differentiable and satisfy $f(0)=f(1)=0$. Prove that
\[
\int_0^1 [f'(x)]^2\,dx \ge 12 \left(\int_0^1 f(x)\,dx\right)^2.
\]
\\
d) Let $f : [0,1] \to \mathbb{R}$ be differentiable and satisfy $f(0)=f(1)=0$. Prove that
\[
\int_0^1 [f'(x)]^2\,dx \ge \pi^2 \int_0^1 [f(x)]^2\,dx.
\]
\\
e) Let $f : [0,1] \to \mathbb{R}$ be twice continuously differentiable and satisfy $f(0)=f'(0)=0$. Prove that
\[
\int_0^1 [f''(x)]^2\,dx \ge 20 \left(\int_0^1 f(x)\,dx\right)^2.
\]
\\
f) Let $f : [0,1] \to \mathbb{R}$ be twice continuously differentiable and satisfy $f(0)=f(1)=0$, $f'(0)=f'(1)=0$. Prove that
\[
\int_0^1 [f''(x)]^2\,dx \ge 720 \left(\int_0^1 f(x)\,dx\right)^2.
\]
\\
g) Let $f : [0,1] \to \mathbb{R}$ be twice continuously differentiable and satisfy $f(0)=f'(0)=0$. Prove that
\[
\int_0^1 [f''(x)]^2\,dx \ge \frac{45}{2} \left(\int_0^1 x f(x)\,dx\right)^2.
\]
\\
h) Let $f : [0,1] \to \mathbb{R}$ be differentiable and satisfy $f(0)=f(1)=0$. Find the smallest constant $C$ such that
\[
\max_{x \in [0,1]} |f(x)| \le C \sqrt{\int_0^1 [f'(t)]^2\,dt}.
\]
\\
i) Let $f : [0,1] \to \mathbb{R}$ be differentiable with $f(0)=f(1)=0$. Find the best constant $C$ such that
\[
|f(c)|^2 \le C \int_0^1 [f'(x)]^2\,dx, \qquad \forall c \in [0,1].
\]
\\
j) Let $f : [0,1] \to \mathbb{R}$ be twice continuously differentiable and satisfy $f(0)=f(1)=0$. Find the best constant $C$ such that
\[
\int_0^1 |f''(x)|\,dx \ge C \max_{x \in [0,1]} |f(x)|.
\]
\\
k) Let $f : [0,1] \to \mathbb{R}$ be twice continuously differentiable, $f(0)=0$, $f'(0)=0$. Find the best constant $C$ such that
\[
f(1)^2 \le C \int_0^1 [f''(x)]^2\,dx.
\]
\\
l) Let $M = \{f \mid f \in C^2[0,1],\ f(0)=f(1)=0,\ f'(0)=1\}$. Determine
\[
\min_{f \in M} \int_0^1 (f''(x))^2\,dx
\]
and the functions for which the minimum is achieved.\\
m) Let $\mathcal{F}$ be the set of continuously differentiable functions $f : [0,1] \to \mathbb{R}$ such that $\displaystyle\int_0^1 f(x)\,dx = 0$, $f \not\equiv 0$. Find the largest constant $K$ such that
\[
\int_0^1 [f'(x)]^2\,dx \ge K \left(\int_0^1 x^2 f(x)\,dx\right)^2,
\]
and determine the case of equality.
\end{problembox}

\stepcounter{probnum}
\begin{problembox}
\textbf{\theprobnum.} Let $M > 0$ be a real number and let $f : [a,b] \to \mathbb{R}$ be a twice differentiable function with continuous second derivative on $[a,b]$, with the property that $|f''(x)| \le M$ for all $x \in [a,b]$. Prove that
\[
\left|\frac{1}{b-a}\int_a^b f(x)\,dx - f\left(\frac{a+b}{2}\right)\right| \le M\,\frac{(b-a)^2}{24}.
\]
\end{problembox}

\stepcounter{probnum}
\begin{problembox}
\textbf{\theprobnum.} Let $a \in [0,1]$. Determine all closed intervals $I \subset [0,1]$ with the property that for every differentiable convex function $f : [0,1] \to \mathbb{R}$ the inequality
\[
\int_0^1 \max\{f(a), f(x)\}\,dx \ge \max_{x \in I} f(x)
\]
holds.
\end{problembox}

\stepcounter{probnum}
\begin{problembox}
\textbf{\theprobnum.} Let $f_0 : [0,1] \to \mathbb{R}$ be a continuous function, and let $f_{n+1}$ be the antiderivative of $f_n$ on $[0,1]$ that vanishes at the origin, for $n \in \mathbb{N}$. If $f_{1982}(1) = \dfrac{1}{1983!}$, prove that there exists $x_0 \in [0,1]$ such that $f_0(x_0) = x_0$.
\end{problembox}

\stepcounter{probnum}
\begin{problembox}
\textbf{\theprobnum.} Let $f : [a,b] \to \mathbb{R}$ be an integrable function. Prove that:\\
1) for each $\alpha \in (0,1)$ there exists at least one $c \in (a,b)$ such that $\displaystyle\alpha \int_a^b f(x)\,dx = \int_a^c f(x)\,dx$;\\
2) if $f$ is continuous on $(a,b)$ and $\displaystyle\int_a^b f(x)\,dx \ne 0$, then for every $n \in \mathbb{N}^*$ there exist $n$ distinct points $x_1, x_2, \ldots, x_n$ in the interval $(a,b)$ such that
\[
\int_a^b f(x)\,dx = \frac{n(b-a)}{\dfrac{1}{f(x_1)} + \dfrac{1}{f(x_2)} + \cdots + \dfrac{1}{f(x_n)}}.
\]
\end{problembox}

\stepcounter{probnum}
\begin{problembox}
\textbf{\theprobnum.} Prove that
\[
\sum_{n=1}^{\infty} \sum_{m=1}^{\infty} \frac{1}{m(m+1)\cdots(m+n)} = \int_{0+0}^{1} \frac{e^x-1}{x}\,dx.
\]
\end{problembox}

\stepcounter{probnum}
\begin{problembox}
\textbf{\theprobnum.} (Generalization of the preceding problem) Prove that for every $m \in \mathbb{N}$ the equality
\[
\sum_{n_m=1}^{\infty} \cdots \sum_{n_1=1}^{\infty} \sum_{n_0=1}^{\infty} \frac{1}{n_0(n_0+n_1)\cdots(n_0+n_1+\cdots+n_m)}
\]
\[
= (-1)^{m-1} \left( \int_{0+0}^{1} \frac{e^x-1}{x}\,dx + \sum_{j=1}^{m-1} \frac{1}{j}\left(1 - e\sum_{i=0}^{j-1} \frac{(-1)^i}{i!}\right) \right)
\]
holds.
\end{problembox}

\stepcounter{probnum}
\begin{problembox}
\textbf{\theprobnum.} Determine the continuous decreasing functions $f : [0,1] \to [0,\infty)$ that simultaneously satisfy the conditions
\[
\int_0^1 f(x)\,dx \le 1 \qquad \text{and} \qquad \int_0^1 f(x^3)\,dx \le 3\left(\int_0^1 x f(x)\,dx\right)^2.
\]
\end{problembox}

\stepcounter{probnum}
\begin{problembox}
\textbf{\theprobnum.} Define $g_0 = 1$ and $g_{n+1}(x) = 1 + \displaystyle\int_0^x g_n(t-t^2)\,dt$ for $x \in [0,1]$.\\
a) Show that $0 \le g_n(x) - g_{n-1}(x) \le \dfrac{x^n}{n!}$.\\
b) Deduce the pointwise convergence of $(g_n)$.\\
c) Prove that the convergence is uniform on $[0,1]$.
\end{problembox}

\stepcounter{probnum}
\begin{problembox}
\textbf{\theprobnum.} Let $f \in C[a,b]$ be such that $\displaystyle\int_a^b x^n f(x)\,dx = 0$ for every $n \in \mathbb{N}$. Show that $f$ is identically zero.
\end{problembox}

\stepcounter{probnum}
\begin{problembox}
\textbf{\theprobnum.} Let $f \in C^1[a,b]$ with $f'(a) \ne 0$. For any $x \in (a,b]$, let $\theta(x) \in [a,x]$ be such that
\[
\int_a^x f(t)\,dt = (x-a) f(\theta(x)).
\]
Calculate $\displaystyle\lim_{x \to a} \frac{\theta(x)-a}{x-a}$.
\end{problembox}

\stepcounter{probnum}
\begin{problembox}
\textbf{\theprobnum.} Let $f : [0,1] \to \mathbb{R}$ be a continuous function such that
\[
\int_0^1 f(x)\,dx = \int_0^1 x f(x)\,dx = \cdots = \int_0^1 x^{n-1} f(x)\,dx = 0
\quad \text{and} \quad \int_0^1 x^n f(x)\,dx = 1.
\]
Show that there exists $a \in [0,1]$ such that $|f(a)| \ge 2^{n(n+1)}$.
\end{problembox}

\stepcounter{probnum}
\begin{problembox}
\textbf{\theprobnum.} Let $f : \mathbb{R} \to \mathbb{R}$ be an additive function (i.e., $f(x+y)=f(x)+f(y)$ for all $x,y \in \mathbb{R}$) that is integrable on every compact interval of $\mathbb{R}$. Show that there exists $a \in \mathbb{R}$ such that $f(x) = ax$ for all $x \in \mathbb{R}$.
\end{problembox}

\stepcounter{probnum}
\begin{problembox}
\textbf{\theprobnum.} Let $f : [0,1] \to \mathbb{R}$ be a continuous function. Show that
\[
\int_0^1 \int_x^1 f(t)\,dt\,dx = \int_0^1 x f(x)\,dx.
\]
\end{problembox}

\stepcounter{probnum}
\begin{problembox}
\textbf{\theprobnum.} Let $f_1, f_2, f_3, f_4 \in \mathbb{R}[x]$. Show that
\[
F(x) = \int_1^x f_1(t)f_2(t)\,dt \int_1^x f_3(t)f_4(t)\,dt - \int_1^x f_1(t)f_4(t)\,dt \int_1^x f_2(t)f_3(t)\,dt
\]
is a polynomial divisible by $(x-1)^4$.
\end{problembox}

\stepcounter{probnum}
\begin{problembox}
\textbf{\theprobnum.} Let $f : [0,\pi] \to \mathbb{R}$ be a continuous function such that
\[
\int_0^{\pi} f(x) \cos x\,dx = \int_0^{\pi} f(x) \sin x\,dx = 0.
\]
Show that the equation $f(x) = 0$ has at least two roots in $(0,\pi)$.
\end{problembox}

\stepcounter{probnum}
\begin{problembox}
\textbf{\theprobnum.} Let $f : [0,1] \to (0,\infty)$ be an increasing function. Prove that
\[
\int_0^1 x f^2(x)\,dx \int_0^1 f(x)\,dx \ge \int_0^1 f^2(x)\,dx \int_0^1 x f(x)\,dx.
\]
\end{problembox}

\stepcounter{probnum}
\begin{problembox}
\textbf{\theprobnum.} Let $f : [-1,1] \to \mathbb{R}$ be a twice differentiable function such that $|f(x)| \le 1$ and $|f''(x)| \le 1$ for all $x \in [-1,1]$. Show that $|f'(x)| \le 2$ for all $x \in [-1,1]$.
\end{problembox}

\stepcounter{probnum}
\begin{problembox}
\textbf{\theprobnum.} Determine the maximum of the expression
\[
\int_0^1 x^2 f(x)\,dx - \int_0^1 x (f(x))^2\,dx
\]
for $f \in C[0,1]$.
\end{problembox}

\stepcounter{probnum}
\begin{problembox}
\textbf{\theprobnum.} Let $I \subseteq \mathbb{R}$ be an open interval with $0 \in I$ and $f : I \to \mathbb{R}$ a continuous function with $f(0)=0$, right-differentiable at $0$.\\
a) Show that $\displaystyle\lim_{n \to \infty} \sum_{k=1}^{n} f\left(\frac{k}{n^2}\right) = \frac{1}{2} f'_d(0)$.\\
b) Special case: calculate $\displaystyle\lim_{n \to \infty} \sum_{k=1}^{n} \left(\sqrt{1+\frac{k}{n^2}} - 1\right)$.
\end{problembox}

\stepcounter{probnum}
\begin{problembox}
\textbf{\theprobnum.} Let $f : (0,\infty) \to \mathbb{R}$ be a differentiable function. Assume there exists $\displaystyle\lim_{x \to \infty} (f(x)+f'(x)) = l \in \mathbb{R}$. Show that
\[
\lim_{x \to \infty} f(x) = l \qquad \text{and} \qquad \lim_{x \to \infty} f'(x) = 0.
\]
\end{problembox}

\stepcounter{probnum}
\begin{problembox}
\textbf{\theprobnum.} Let $f : \mathbb{R} \to \mathbb{R}$ be a twice continuously differentiable function. If $f(x+y)f(x-y) \le f^2(x)$ for all $x,y \in \mathbb{R}$, show that $f(x)f''(x) \le (f'(x))^2$ for all $x \in \mathbb{R}$.
\end{problembox}

\stepcounter{probnum}
\begin{problembox}
\textbf{\theprobnum.} Let $I \subseteq \mathbb{R}$ be an interval and $f : I \to \mathbb{R}$ an arbitrary function. Show that $f$ has at most a countable set of strict extrema.
\end{problembox}

\stepcounter{probnum}
\begin{problembox}
\textbf{\theprobnum.} Suppose the differentiable functions $a,b,f,g : \mathbb{R} \to \mathbb{R}$ satisfy $f(x) \ge 0$, $f'(x) \ge 0$, $g(x) > 0$, $g'(x) > 0$ for all $x \in \mathbb{R}$. Let $\displaystyle\lim_{x \to \infty} a(x) = A > 0$, $\displaystyle\lim_{x \to \infty} b(x) = B > 0$, $\displaystyle\lim_{x \to \infty} f(x) = \lim_{x \to \infty} g(x) = \infty$, and assume
\[
\lim_{x \to \infty} \left( \frac{f'(x)}{g'(x)} + a(x) \frac{f(x)}{g(x)} - b(x) \right) = 0.
\]
Show that $\displaystyle\lim_{x \to \infty} \frac{f(x)}{g(x)} = \frac{B}{A+1}$.
\end{problembox}

\stepcounter{probnum}
\begin{problembox}
\textbf{\theprobnum.} Let $k > 1$ be a real number. Calculate the limits:\\
a) $L = \displaystyle\lim_{n \to \infty} \int_0^1 \sqrt[n]{kx+k-1}\,dx$;\\
b) $\displaystyle\lim_{n \to \infty} n \left(L - \int_0^1 \sqrt[n]{kx+k-1}\,dx\right)$.
\end{problembox}

\stepcounter{probnum}
\begin{problembox}
\textbf{\theprobnum.} Let $f : [0,1] \to \mathbb{R}$ be the function $f(x) = (x^2+1)e^x$. Calculate
\[
\lim_{n \to \infty} \int_0^1 \left(f\left(\frac{x}{n}\right)\right)^n dx.
\]
\end{problembox}

\stepcounter{probnum}
\begin{problembox}
\textbf{\theprobnum.} Determine
\[
\lim_{n \to \infty} \left(\int_0^1 e^{x^2/n}\,dx\right)^n.
\]
\end{problembox}

\stepcounter{probnum}
\begin{problembox}
\textbf{\theprobnum.} Assume $f(x) \ge 0$ and $\displaystyle\int_0^{\infty} f^2(x)\,dx < \infty$. Define $F(x) = \displaystyle\int_0^x f(t)\,dt$. Prove that
\[
\int_0^{\infty} \left(\frac{F(x)}{x}\right)^2 dx \le 4 \int_0^{\infty} f^2(x)\,dx.
\]
\end{problembox}

\stepcounter{probnum}
\begin{problembox}
\textbf{\theprobnum.} Let $f : \mathbb{R} \to \mathbb{R}$ be continuously differentiable, satisfying $f(0)=0$, $f(1)=1$, and $f'(x) > 0$. Define
\[
L = \int_0^1 \sqrt{1+[f'(x)]^2}\,dx.
\]
Prove that
\[
\int_0^1 x \sqrt{1+[f'(x)]^2}\,dx \le \frac{L}{2} + \frac{1}{4}.
\]
\end{problembox}

\stepcounter{probnum}
\begin{problembox}
\textbf{\theprobnum.} (Gronwall's inequality) Let $u(t)$ be a nonnegative continuous function on $[a,b]$. Assume there exist constants $C \ge 0$, $K \ge 0$ such that
\[
u(t) \le C + K \int_a^t u(s)\,ds, \qquad \forall t \in [a,b].
\]
Prove that $u(t) \le C e^{K(t-a)}$.
\end{problembox}

\stepcounter{probnum}
\begin{problembox}
\textbf{\theprobnum.} Let $f : [a,b] \to \mathbb{R}$ be continuous. Prove that there exists $c \in (a,b)$ such that
\[
f(c) \left(\int_c^b f(x)\,dx - \int_a^c f(x)\,dx\right) = 0.
\]
\end{problembox}

\stepcounter{probnum}
\begin{problembox}
\textbf{\theprobnum.} Let $f : [0,1] \to \mathbb{R}$ be continuous and positive. Prove that there exists $c \in (0,1)$ such that
\[
\int_0^c f(x)\,dx = (1-c) f(c).
\]
\end{problembox}

\stepcounter{probnum}
\begin{problembox}
\textbf{\theprobnum.} Let $f : [0,1] \to \mathbb{R}$ be continuous and satisfy $\displaystyle\int_0^1 f(x)\,dx = 0$. Prove that there exists $c \in (0,1)$ such that
\[
f(c) = \int_0^c f(x)\,dx.
\]
\end{problembox}

\stepcounter{probnum}
\begin{problembox}
\textbf{\theprobnum.} Let $f : [0,1] \to \mathbb{R}$ be continuous and assume $f(1)=0$. Prove that there exists at least one $c \in (0,1)$ such that
\[
f(c) = \int_0^c f(x)\,dx.
\]
\end{problembox}

\stepcounter{probnum}
\begin{problembox}
\textbf{\theprobnum.} Let $f : [0,1] \to \mathbb{R}$ be continuous. Prove that for every positive integer $n$, there exists $c \in (0,1)$ such that
\[
\int_0^c x^{n-1} f(x)\,dx = \frac{c^n}{n} f(c).
\]
\end{problembox}

\stepcounter{probnum}
\begin{problembox}
\textbf{\theprobnum.} Let $f : [0,\tfrac{\pi}{2}] \to \mathbb{R}$ be continuous and satisfy
\[
\int_0^{\pi/2} f(x) \sin x\,dx = \int_0^{\pi/2} f(x) \cos x\,dx.
\]
Prove that there exists $c \in (0,\tfrac{\pi}{2})$ such that
\[
f(c) = \left(\int_0^c f(x)\,dx\right)(\sin c + \cos c).
\]
\end{problembox}

\stepcounter{probnum}
\begin{problembox}
\textbf{\theprobnum.} Let $f : [0,1] \to \mathbb{R}$ be differentiable and $k \ne 0$. Assume $f(0)=0$ and
\[
f(1) = k \int_0^1 e^{k(1-x)} f(x)\,dx.
\]
Prove that there exists $c \in (0,1)$ such that $f'(c) = k f(c)$.
\end{problembox}

\stepcounter{probnum}
\begin{problembox}
\textbf{\theprobnum.} Let $f : [a,b] \to \mathbb{R}$ be continuous and assume $\displaystyle\int_a^b f(x)\,dx \ne 0$. Prove that there exists $c \in (a,b)$ such that
\[
f(c) \left[\frac{1}{\int_a^c f(x)\,dx} - \frac{1}{\int_c^b f(x)\,dx}\right] = \frac{1}{c-a} - \frac{1}{b-c}.
\]
\end{problembox}

\stepcounter{probnum}
\begin{problembox}
\textbf{\theprobnum.} Let $f : [0,1] \to \mathbb{R}$ be continuous. Prove that there exists $c \in (0,1)$ such that
\[
2c^2 f(c) \left(\int_0^c f(x)\,dx - \int_c^1 f(x)\,dx\right) = \int_0^c f(x)\,dx \int_c^1 f(x)\,dx.
\]
\end{problembox}

\stepcounter{probnum}
\begin{problembox}
\textbf{\theprobnum.} Let $f : [0,1] \to \mathbb{R}$ be continuously differentiable, $f(0)=0$, $f(1)=1$, and $f'(x)+f(x) > 0$. Find the minimum value of
\[
I = \int_0^1 \frac{1}{f'(x)+f(x)}\,dx.
\]
\end{problembox}

\stepcounter{probnum}
\begin{problembox}
\textbf{\theprobnum.} Under the same hypotheses as the preceding problem, find the minimum value of
\[
J = \int_0^1 \frac{1}{[f'(x)+f(x)]^2}\,dx.
\]
\end{problembox}

\stepcounter{probnum}
\begin{problembox}
\textbf{\theprobnum.} Let $f$ be continuous on $[0,1]$ and satisfy $1 \le f(x) \le 3$. Prove that
\[
\int_0^1 f(x)\,dx \int_0^1 \frac{1}{f(x)}\,dx \le \frac{4}{3}.
\]
\end{problembox}

\stepcounter{probnum}
\begin{problembox}
\textbf{\theprobnum.} (Polya--Szeg\H{o} inequality) Let $f,g$ be positive integrable functions on $[0,1]$ and assume
\[
0 < m \le \frac{f(x)}{g(x)} \le M.
\]
Prove that
\[
\int_0^1 f^2(x)\,dx \int_0^1 g^2(x)\,dx \le \frac{(m+M)^2}{4mM} \left(\int_0^1 f(x)g(x)\,dx\right)^2.
\]
\end{problembox}

\stepcounter{probnum}
\begin{problembox}
\textbf{\theprobnum.} Let $f$ be continuous on $[0,1]$ and satisfy $1 \le f(x) \le 3$. Find the maximum value of
\[
P = \int_0^1 f(x)\,dx + \int_0^1 \frac{1}{f(x)}\,dx - \int_0^1 f(x)\,dx \int_0^1 \frac{1}{f(x)}\,dx.
\]
\end{problembox}

\stepcounter{probnum}
\begin{problembox}
\textbf{\theprobnum.} Let $f : [0,1] \to \mathbb{R}$ be continuously differentiable and satisfy $1 \le f(x) \le 3$. Prove that
\[
\int_0^1 f(x)\,dx \int_0^1 \frac{1}{f(x)}\,dx \le \min\left\{\frac{4}{3},\ 1 + \frac{1}{\pi^2}\int_0^1 \left(\frac{f'(x)}{f(x)}\right)^2 dx\right\}.
\]
\end{problembox}

\stepcounter{probnum}
\begin{problembox}
\textbf{\theprobnum.} Let $f, g : \mathbb{R} \to \mathbb{R}$ be nonconstant, differentiable functions satisfying
\[
f(x+y) = f(x)f(y) - g(x)g(y), \qquad g(x+y) = f(x)g(y) + g(x)f(y),
\]
for all $x, y \in \mathbb{R}$, and $f'(0) = 0$. Prove that $(f(x))^2 + (g(x))^2 = 1$ for all $x \in \mathbb{R}$.
\end{problembox}

\stepcounter{probnum}
\begin{problembox}
\textbf{\theprobnum.} Let $f$ be a real function on the real line with a continuous third derivative. Prove that there exists a point $a \in \mathbb{R}$ such that
\[
f(a) \cdot f'(a) \cdot f''(a) \cdot f'''(a) \ge 0.
\]
\end{problembox}

\stepcounter{probnum}
\begin{problembox}
\textbf{\theprobnum.} Suppose that $f : [0,1] \to \mathbb{R}$ has a continuous derivative and that $\displaystyle\int_0^1 f(x)\,dx = 0$. Prove that for every $\alpha \in (0,1)$,
\[
\left|\int_0^{\alpha} f(x)\,dx\right| \le \frac{1}{8} \max_{0 \le x \le 1} |f'(x)|.
\]
\end{problembox}

\stepcounter{probnum}
\begin{problembox}
\textbf{\theprobnum.} Let $f : [0,1] \to \mathbb{R}$ be a continuous function satisfying $x f(y) + y f(x) \le 1$ for every $x, y \in [0,1]$.\\
a) Show that $\displaystyle\int_0^1 f(x)\,dx \le \frac{\pi}{4}$.\\
b) Find such a function for which equality occurs.
\end{problembox}

\stepcounter{probnum}
\begin{problembox}
\textbf{\theprobnum.} Let $f : \mathbb{R} \to \mathbb{R}$ be a continuous function. A point $x$ is called a \emph{shadow point} if there exists a point $y \in \mathbb{R}$ with $y > x$ such that $f(y) > f(x)$. Let $a < b$ be real numbers and suppose that all the points of the open interval $I = (a,b)$ are shadow points, while $a$ and $b$ are not shadow points. Prove that:\\
a) $f(x) \le f(b)$ for all $a < x < b$;\\
b) $f(a) = f(b)$.
\end{problembox}

\stepcounter{probnum}
\begin{problembox}
\textbf{\theprobnum.} Let $(a_n) \subset \left[\dfrac{1}{2}, 1\right]$. Define the sequence $x_0 = 0$, and
\[
x_{n+1} = \frac{a_{n+1} + x_n}{1 + a_{n+1} x_n}.
\]
Is this sequence convergent? If yes, find its limit.
\end{problembox}

\stepcounter{probnum}
\begin{problembox}
\textbf{\theprobnum.} Let $f : \mathbb{R} \to \mathbb{R}$ be a function such that $|f(x) - f(y)| \le |x-y|$ for all $x, y \in \mathbb{R}$. Prove that if for any real $x$ the sequence $x, f(x), f(f(x)), \ldots$ is an arithmetic progression, then there is $a \in \mathbb{R}$ such that $f(x) = x + a$ for all $x \in \mathbb{R}$.
\end{problembox}

\stepcounter{probnum}
\begin{problembox}
\textbf{\theprobnum.} Let $f : [-1,1] \to \mathbb{R}$ be a continuous function having finite derivative at $0$, and let
\[
I(h) = \int_{-h}^{h} f(x)\,dx, \qquad h \in [0,1].
\]
Prove that:\\
a) There exists $M > 0$ such that $|I(h) - 2f(0)h| \le Mh^2$ for any $h \in [0,1]$.\\
b) The sequence $(a_n)_{n \ge 1}$ defined by $\displaystyle a_n = \sum_{k=1}^{n} \sqrt{k}\, \left|I\!\left(\frac{1}{k}\right)\right|$ is convergent if and only if $f(0) = 0$.
\end{problembox}

\stepcounter{probnum}
\begin{problembox}
\textbf{\theprobnum.} Let $f : \mathbb{R} \to \mathbb{R}$ be a function, two times derivable on $\mathbb{R}$, for which there exists $c \in \mathbb{R}$ such that
\[
\frac{f(b)-f(a)}{b-a} = f'(c)
\]
for all $a \ne b \in \mathbb{R}$. Prove that $f''(c) = 0$.
\end{problembox}

\stepcounter{probnum}
\begin{problembox}
\textbf{\theprobnum.} a) Consider a differentiable and convex function $f : [0,\infty) \to [0,\infty)$. Show that if $f(x) \le x$ for every $x \ge 0$, then $f'(x) \le 1$ for every $x \ge 0$.\\
b) Determine differentiable and convex functions $f : [0,\infty) \to [0,\infty)$ which have the property that $f(0) = 0$ and $f'(x) f(f(x)) = x$ for every $x \ge 0$.
\end{problembox}

\stepcounter{probnum}
\begin{problembox}
\textbf{\theprobnum.} Let $f : \mathbb{R} \to [0,\infty)$ be a twice differentiable function whose second derivative is continuous, such that
\[
\lim_{x \to \infty} \left[f(x) + f'(x) + f''(x)\right] = C.
\]
Prove that $\displaystyle\lim_{x \to \infty} f(x) = C$.
\end{problembox}

\stepcounter{probnum}
\begin{problembox}
\textbf{\theprobnum.} Let $f$ be a real function with a continuous third derivative such that $f(x)$, $f'(x)$, $f''(x)$, and $f'''(x)$ are positive for all $x$. Suppose that $f'''(x) \le f(x)$ for all $x$. Show that $f'(x) < 2f(x)$ for all $x$.
\end{problembox}

\stepcounter{probnum}
\begin{problembox}
\textbf{\theprobnum.} Suppose that $f$ is a function on the interval $[1,3]$ such that $-1 \le f(x) \le 1$ for all $x$ and $\displaystyle\int_1^3 f(x)\,dx = 0$. How large can $\displaystyle\int_1^3 \frac{f(x)}{x}\,dx$ be?
\end{problembox}

\stepcounter{probnum}
\begin{problembox}
\textbf{\theprobnum.} a) Prove that there exists a differentiable function $f : (0,\infty) \to (0,\infty)$ such that $f(f'(x)) = x$ for all $x > 0$.\\
b) Prove that there is no differentiable function $f : \mathbb{R} \to \mathbb{R}$ such that $f(f'(x)) = x$ for all $x \in \mathbb{R}$.
\end{problembox}

\stepcounter{probnum}
\begin{problembox}
\textbf{\theprobnum.} Let $a > 0$ and let $\mathcal{F} = \{f : [0,1] \to \mathbb{R} \mid f \text{ is concave and } f(0) = 1\}$. Determine
\[
\min_{f \in \mathcal{F}} \left\{ \left(\int_0^1 f(x)\,dx\right)^2 - (a+1) \int_0^1 x^{2a} f(x)\,dx \right\}.
\]
\end{problembox}

\stepcounter{probnum}
\begin{problembox}
\textbf{\theprobnum.} Let $f : \mathbb{R} \to \mathbb{R}$ be a function such that there exists $M > 0$ satisfying
\[
|f(x_1 + \cdots + x_n) - f(x_1) - \cdots - f(x_n)| \le M, \qquad \forall x_1, \ldots, x_n \in \mathbb{R}.
\]
Prove that $f(x+y) = f(x) + f(y)$ for all $x, y \in \mathbb{R}$.
\end{problembox}

\stepcounter{probnum}
\begin{problembox}
\textbf{\theprobnum.} Let $f : (0,1) \to \mathbb{R}$ be a function satisfying $\displaystyle\lim_{x \to 0^+} f(x) = 0$. If there exists $\lambda \in (0,1)$ such that $\displaystyle\lim_{x \to 0^+} \frac{|f(x) - f(\lambda x)|}{x} = 0$, prove that
\[
\lim_{x \to 0^+} \frac{f(x)}{x} = 0.
\]
\end{problembox}

\stepcounter{probnum}
\begin{problembox}
\textbf{\theprobnum.} Determine twice differentiable functions $f : \mathbb{R} \to \mathbb{R}$ which verify the relation
\[
(f'(x))^2 + f''(x) \le 0, \qquad \forall x \in \mathbb{R}.
\]
\end{problembox}

\stepcounter{probnum}
\begin{problembox}
\textbf{\theprobnum.} For an arbitrary number $x_0 \in (0,\pi)$, define recursively the sequence $(x_n)_{n \ge 0}$ by $x_{n+1} = \sin(x_n)$. Compute
\[
\lim_{n \to \infty} \sqrt{n}\, x_n.
\]
\end{problembox}

\stepcounter{probnum}
\begin{problembox}
\textbf{\theprobnum.} Let $a$ be a real number, and $f : \mathbb{R} \to \mathbb{R}$ a non-decreasing function. Prove that $f$ is continuous at $a$ if and only if there exists a sequence $(a_n)_{n \ge 1}$ of positive real numbers such that
\[
\int_a^{a+a_n} f(x)\,dx + \int_a^{a-a_n} f(x)\,dx \le \frac{a_n}{n}
\]
for all natural numbers $n$.
\end{problembox}

\stepcounter{probnum}
\begin{problembox}
\textbf{\theprobnum.} Let $f : \mathbb{R} \to \mathbb{R}$ be a surjective function that has the property that if the sequence $(f(x_n))_{n \ge 1}$ is convergent, then the sequence $(x_n)_{n \ge 1}$ is convergent. Prove that $f$ is continuous.
\end{problembox}

\stepcounter{probnum}
\begin{problembox}
\textbf{\theprobnum.} Let $f : [0,\infty) \to (0,\infty)$ be a continuous function.\\
a) Show that there exists a natural number $n_0$ and a sequence of positive real numbers $(x_n)_{n \ge n_0}$ satisfying
\[
n \int_0^{x_n} f(t)\,dt = 1, \qquad \forall n > n_0.
\]
b) Prove that the sequence $(n x_n)_{n \ge n_0}$ is convergent and find its limit.
\end{problembox}

\stepcounter{probnum}
\begin{problembox}
\textbf{\theprobnum.} Let $f : [a,b] \to \mathbb{R}$ be a function with the Intermediate Value Property such that $f(a) \ast f(b) < 0$. Show that there exist $\alpha, \beta$ such that $a < \alpha < \beta < b$ and
\[
f(\alpha) + f(\beta) = f(\alpha) \ast f(\beta).
\]
\end{problembox}

\stepcounter{probnum}
\begin{problembox}
\textbf{\theprobnum.} Let $f : [0,1] \to (0,1)$ be a surjective function.\\
a) Prove that $f$ has at least one point of discontinuity.\\
b) Given that $f$ admits a limit at any point of the interval $[0,1]$, show that it has at least two points of discontinuity.
\end{problembox}

\stepcounter{probnum}
\begin{problembox}
\textbf{\theprobnum.} Let $f : [0,\infty) \to (0,\infty)$ be an increasing function and $g : [0,\infty) \to \mathbb{R}$ a twice differentiable function such that $g''$ is continuous and
\[
g''(x) + f(x) g(x) = 0, \qquad \forall x \ge 0.
\]
a) Provide an example of such functions, with $g \not\equiv 0$.\\
b) Prove that $g$ is bounded.
\end{problembox}

\stepcounter{probnum}
\begin{problembox}
\textbf{\theprobnum.} Find the number of trapeziums that can be formed using the vertices of a regular $n$-sided polygon.
\end{problembox}

\stepcounter{probnum}
\begin{problembox}
\textbf{\theprobnum.} Let $f : [0,\infty) \to \mathbb{R}$ be a continuous function, constant on $\mathbb{Z}_{\ge 0}$. For any $0 \le a < b < c < d$ which satisfy $f(a) = f(c)$ and $f(b) = f(d)$, we also have
\[
f\left(\frac{a+b}{2}\right) = f\left(\frac{c+d}{2}\right).
\]
Prove that $f$ is a constant function.
\end{problembox}

\stepcounter{probnum}
\begin{problembox}
\textbf{\theprobnum.} Let $f_n$ be a sequence of real-valued $C^1$ functions on $[0,1]$ such that for all $n \in \mathbb{N}$ the following hold:
\[
|f_n'(x)| \le \frac{1}{\sqrt{x}} \quad (0 < x \le 1), \qquad \int_0^1 f_n(x)\,dx = 0.
\]
Prove that $f_n$ has a convergent subsequence that converges uniformly on $[0,1]$.
\end{problembox}

\stepcounter{probnum}
\begin{problembox}
\textbf{\theprobnum.} Let $\chi_{\mathbb{Q}}$ denote the characteristic function of the rationals in $[0,1]$. Does there exist a sequence of continuous functions $f_n : [0,1] \to \mathbb{R}$ such that $f_n$ converges to $\chi_{\mathbb{Q}}$ pointwise?
\end{problembox}

\stepcounter{probnum}
\begin{problembox}
\textbf{\theprobnum.} Let $f : [0,1] \to \mathbb{R}$ be a continuous function such that
\[
\int_0^1 f(t)\,dt = \int_0^1 t f(t)\,dt.
\]
Prove that there exists $c \in (0,1)$ such that
\[
\int_0^c f(t)\,dt = \frac{c}{2} f(c).
\]
\end{problembox}

\stepcounter{probnum}
\begin{problembox}
\textbf{\theprobnum.} Let $f : [0,1] \to \mathbb{R}$ be a continuous function such that $\displaystyle\int_0^1 f(t)\,dt = \int_0^1 t f(t)\,dt$. Prove that there exists $c \in (0,1)$ such that
\[
c f(c) = 2 \int_c^1 f(t)\,dt.
\]
\end{problembox}

\stepcounter{probnum}
\begin{problembox}
\textbf{\theprobnum.} Let $f : \mathbb{R} \to \mathbb{R}$ be a differentiable function such that
\[
f'(x) = f^2(x) f(-x).
\]
Find an explicit formula for $f$.
\end{problembox}

\stepcounter{probnum}
\begin{problembox}
\textbf{\theprobnum.} Let $f : [0,1] \to \mathbb{R}$ be a continuous function such that $\displaystyle\int_0^1 f(t)\,dt = 1$ and
\[
\int_0^1 (1-f(x)) e^{-f(x)}\,dx \le 0.
\]
Prove that $f(x) = 1$ for all $x \in [0,1]$.
\end{problembox}

\stepcounter{probnum}
\begin{problembox}
\textbf{\theprobnum.} Let $f : [a,b] \to [0,+\infty)$ be a continuous and not everywhere $0$ function. Prove that
\[
\lim_{n \to +\infty} \frac{\displaystyle\int_a^b f^{n+1}(t)\,dt}{\displaystyle\int_a^b f^n(t)\,dt} = \sup_{x \in [a,b]} f(x).
\]
\end{problembox}

\stepcounter{probnum}
\begin{problembox}
\textbf{\theprobnum.} Examine if there exists a continuous function $f : [1,+\infty) \to \mathbb{R}$ such that $f(x) > 0$ for all $x \in [1,+\infty)$ and $\displaystyle\int_1^{\infty} f(t)\,dt$ converges whereas $\displaystyle\int_1^{\infty} f^2(t)\,dt$ diverges.
\end{problembox}

\stepcounter{probnum}
\begin{problembox}
\textbf{\theprobnum.} Let $0 < x_1 < 1$ and $x_{n+1} = x_n(1-x_n)$ for $n = 1,2,3,\ldots$. Show that
\[
\lim_{n \to \infty} n x_n = 1.
\]
\end{problembox}

\stepcounter{probnum}
\begin{problembox}
\textbf{\theprobnum.} a) Let $(x_n)_{n \ge 1}$ be a sequence such that $x_1 \in (0,1)$ and $x_{n+1} = x_n(1-x_n^2)$ for $n \ge 1$. Evaluate
\[
\lim_{n \to \infty} \sqrt{n}\, x_n.
\]
b) Let $(x_n)_{n \ge 1}$ be a sequence such that $x_1 = x > 1$ and $x_{n+1} = x_n + \sqrt{x_n - 1}$ for $n \ge 1$. Evaluate
\[
\lim_{n \to \infty} \frac{4x_n - n^2}{n \log n}.
\]
\end{problembox}

\stepcounter{probnum}
\begin{problembox}
\textbf{\theprobnum.} Consider the harmonic sequence $(H_n)_{n \ge 1}$, $H_n = 1 + \dfrac{1}{2} + \cdots + \dfrac{1}{n}$, $n \ge 1$. Show that the sequence $x_n = \{H_n\}$ diverges, where $\{x\}$ denotes the fractional part of the real number $x$.
\end{problembox}

\stepcounter{probnum}
\begin{problembox}
\textbf{\theprobnum.} Let $(a_n)_{n \ge 1}$ be a decreasing sequence of positive reals such that $\displaystyle\sum_{n=1}^{\infty} a_n$ converges. Show that $\displaystyle\lim_{n \to \infty} n a_n = 0$.
\end{problembox}

\stepcounter{probnum}
\begin{problembox}
\textbf{\theprobnum.} Evaluate the following limits:\\
a) $\displaystyle\lim_{n \to \infty} \left(\frac{\left(1+\frac{1}{n}\right)^n}{e}\right)^n$;\\
b) $\displaystyle\lim_{n \to \infty} \prod_{k=1}^{n} \left(\frac{k}{k+n}\right)^{e^{1999/n}-1}$.
\end{problembox}

\stepcounter{probnum}
\begin{problembox}
\textbf{\theprobnum.} Define the sequence $x_1, x_2, \ldots$ inductively by $x_1 = \sqrt{5}$ and $x_{n+1} = x_n^2 - 2$ for each $n \ge 1$. Compute
\[
\lim_{n \to \infty} \frac{x_1 x_2 \cdots x_n}{x_{n+1}}.
\]
\end{problembox}

\stepcounter{probnum}
\begin{problembox}
\textbf{\theprobnum.} Suppose that $a_0 = 1$ and that $a_{n+1} = a_n + e^{-a_n}$ for $n = 0, 1, 2, \ldots$. Does the sequence $(a_n - \log n)_{n \ge 1}$ have a finite limit?
\end{problembox}

\stepcounter{probnum}
\begin{problembox}
\textbf{\theprobnum.} Let $k$ be an integer greater than $1$. Suppose that $a_0 > 0$, and define the sequence $(a_n)_{n \ge 0}$ by
\[
a_{n+1} = a_n + \frac{1}{\sqrt[k]{a_n}}, \qquad n \ge 0.
\]
Evaluate
\[
\lim_{n \to \infty} \frac{a_n^{k+1}}{n^k}.
\]
\end{problembox}

\stepcounter{probnum}
\begin{problembox}
\textbf{\theprobnum.} a) Does there exist a sequence $(a_n)_{n \ge 1}$ of positive reals such that the series $\displaystyle\sum_{n=1}^{\infty} \frac{1}{a_n}$ converges, and $\displaystyle\prod_{k=1}^{n} a_k \le n^n$ for all $n \ge 1$?\\
b) Does there exist a sequence $(a_n)_{n \ge 1}$ of positive reals such that $\displaystyle\sum_{k=1}^{n} \frac{1}{a_k} \le 2008$, and $\displaystyle\sum_{k=1}^{n} a_k \le n^2$ for all $n \ge 1$?
\end{problembox}

\stepcounter{probnum}
\begin{problembox}
\textbf{\theprobnum.} Let $(a_n)_{n \ge 1}$ be a sequence of positive reals such that the series $\displaystyle\sum_{n=1}^{\infty} a_n$ converges. Show that the series $\displaystyle\sum_{n=1}^{\infty} a_n^{n/(n+1)}$ converges also.
\end{problembox}

\stepcounter{probnum}
\begin{problembox}
\textbf{\theprobnum.} Let $(a_n)_{n \ge 1}$ be a decreasing sequence of positive reals. Let $s_n = a_1 + a_2 + \cdots + a_n$ and
\[
b_n = \frac{1}{a_{n+1}} - \frac{1}{a_n}, \qquad n \ge 1.
\]
Prove that if $(s_n)_{n \ge 1}$ is convergent, then $(b_n)_{n \ge 1}$ is unbounded.
\end{problembox}

\stepcounter{probnum}
\begin{problembox}
\textbf{\theprobnum.} Let $(x_n)_{n \ge 2}$ be a sequence of real numbers such that $x_2 > 0$ and
\[
x_{n+1} = -1 + \sqrt[n]{1+n x_n}, \qquad n \ge 2.
\]
Show that:\\
a) $\displaystyle\lim_{n \to \infty} x_n = 0$;\\
b) $\displaystyle\lim_{n \to \infty} \frac{x_{n+1}}{x_n} = 1$ and $\displaystyle\lim_{n \to \infty} n x_n = 0$.
\end{problembox}

\stepcounter{probnum}
\begin{problembox}
\textbf{\theprobnum.} Let $(a_n)_{n \ge 1}$ be a sequence of positive reals. Show that
\[
\limsup_{n \to \infty} n\left(\frac{1+a_{n+1}}{a_n} - 1\right) \ge 1,
\]
and
\[
\limsup_{n \to \infty} \left(\frac{a_1 + a_{n+1}}{a_n}\right)^n \ge e.
\]
\end{problembox}

\stepcounter{probnum}
\begin{problembox}
\textbf{\theprobnum.} Let $(a_n)_{n \ge 1}$ be a sequence such that $x_n = \displaystyle\sum_{k=1}^{n} a_k^2$ converges, while $y_n = \displaystyle\sum_{k=1}^{n} a_k$ is unbounded. Show that the sequence $(b_n)_{n \ge 1}$, $b_n = \{y_n\}$, diverges, where $\{x\}$ denotes the fractional part of the real number $x$.
\end{problembox}

\stepcounter{probnum}
\begin{problembox}
\textbf{\theprobnum.} Show that if the series $\displaystyle\sum_{n=1}^{\infty} \frac{1}{p_n}$ is convergent, where $p_1, p_2, \ldots, p_n$ are positive real numbers, then the series
\[
\sum_{n=1}^{\infty} \frac{n^2}{(p_1+p_2+\cdots+p_n)^2} \cdot p_n
\]
is also convergent.
\end{problembox}

\stepcounter{probnum}
\begin{problembox}
\textbf{\theprobnum.} Determine, with proof, whether the series $\displaystyle\sum_{n=1}^{\infty} \frac{1}{n^{1.8+\sin n}}$ converges or diverges.
\end{problembox}

\stepcounter{probnum}
\begin{problembox}
\textbf{\theprobnum.} (Carleman's inequality) Let $a_1, a_2, \ldots, a_n$ be a sequence of nonnegative real numbers. Prove that
\[
\sum_{n=1}^{\infty} \sqrt[n]{a_1 a_2 \cdots a_n} \le e \sum_{n=1}^{\infty} a_n.
\]
\end{problembox}

\stepcounter{probnum}
\begin{problembox}
\textbf{\theprobnum.} (Hardy's inequality) Let $a_1, a_2, \ldots, a_n$ be a sequence of nonnegative real numbers and $p > 1$. Prove that
\[
\sum_{n=1}^{\infty} \left(\frac{a_1+a_2+\cdots+a_n}{n}\right)^p \le \left(\frac{p}{p-1}\right)^p \sum_{n=1}^{\infty} a_n^p.
\]
\end{problembox}

\stepcounter{probnum}
\begin{problembox}
\textbf{\theprobnum.} Let $f : [0,\infty) \to \mathbb{R}$ be a function satisfying $f(x) e^{f(x)} = x$ for $x \ge 0$. Prove that:\\
a) $f$ is monotone;\\
b) $\displaystyle\lim_{x \to \infty} f(x) = \infty$;\\
c) $\dfrac{f(x)}{\log x}$ tends to $1$ as $x \to \infty$.
\end{problembox}

\stepcounter{probnum}
\begin{problembox}
\textbf{\theprobnum.} Let $f : (0,\infty) \to \mathbb{R}$ be a function satisfying
\[
\lim_{x \to \infty} \frac{f(x)}{x^k} = a, \qquad k \ne 1,
\]
and
\[
\lim_{x \to \infty} \frac{f(x+1)-f(x)}{x^{k-1}} = b \in \mathbb{R}.
\]
Prove that $b = ka$.
\end{problembox}

\stepcounter{probnum}
\begin{problembox}
\textbf{\theprobnum.} Let $f : (a,b) \to \mathbb{R}$ be a function such that $\displaystyle\lim_{x \to x_0} f(x)$ exists for any $x_0 \in [a,b]$. Prove that $f$ is bounded if and only if for all $x_0 \in [a,b]$, $\displaystyle\lim_{x \to x_0} f(x)$ is finite.
\end{problembox}

\stepcounter{probnum}
\begin{problembox}
\textbf{\theprobnum.} Let $a, b \in \left(0, \dfrac{1}{2}\right)$ and $f : \mathbb{R} \to \mathbb{R}$ be a continuous function such that
\[
f(f(x)) = a f(x) + bx,
\]
for all $x$. Prove that $f(0) = 0$ and find all such functions.
\end{problembox}

\stepcounter{probnum}
\begin{problembox}
\textbf{\theprobnum.} Prove that there is no continuous function $f : \mathbb{R} \to \mathbb{R}$ such that $f(x) \in \mathbb{Q} \implies f(x+1) \in \mathbb{R} \setminus \mathbb{Q}$.
\end{problembox}

\stepcounter{probnum}
\begin{problembox}
\textbf{\theprobnum.} Let $f : \mathbb{R} \to \mathbb{R}$ be a continuous function such that
\[
f(x) \le f\left(x + \frac{1}{n}\right),
\]
for every real $x$ and positive integer $n$. Show that $f$ is nondecreasing.
\end{problembox}

\stepcounter{probnum}
\begin{problembox}
\textbf{\theprobnum.} Find all functions $f : (0,\infty) \to (0,\infty)$ subject to the conditions:\\
(i) $f(f(f(x))) + 2x = f(3x)$ for all $x > 0$;\\
(ii) $\displaystyle\lim_{x \to \infty} (f(x) - x) = 0$.
\end{problembox}

\stepcounter{probnum}
\begin{problembox}
\textbf{\theprobnum.} Let $f(x)$ be a continuous function such that $f(2x^2-1) = 2xf(x)$ for all $x$. Show that $f(x) = 0$ for $-1 \le x \le 1$.
\end{problembox}

\stepcounter{probnum}
\begin{problembox}
\textbf{\theprobnum.} Calculate
\[
\sum_{n=1}^{\infty} \ln\left(1+\frac{1}{n}\right) \ln\left(1+\frac{1}{2n}\right) \ln\left(1+\frac{1}{2n+1}\right).
\]
\end{problembox}

\stepcounter{probnum}
\begin{problembox}
\textbf{\theprobnum.} Let $f(x) = \displaystyle\sum_{k=1}^{n} a_k \sin kx$, with $a_1, a_2, \ldots, a_n$ real numbers and $n \ge 1$. Prove that if $f(x) \le |\sin x|$ for all $x \in \mathbb{R}$, then
\[
\left|\sum_{k=1}^{n} k a_k\right| \le 1.
\]
\end{problembox}

\stepcounter{probnum}
\begin{problembox}
\textbf{\theprobnum.} Let $f : [a,b] \to \mathbb{R}$ be a function, continuous on $[a,b]$ and twice differentiable on $(a,b)$. If $f(a) = f(b)$ and $f'(a) = f'(b)$, prove that for every real number $\lambda$, the equation
\[
f''(x) - \lambda (f'(x))^2 = 0
\]
has at least one zero in the interval $(a,b)$.
\end{problembox}

\stepcounter{probnum}
\begin{problembox}
\textbf{\theprobnum.} Let $\alpha > 1$ be a real number, and let $(u_n)_{n \ge 1}$ be a sequence of positive numbers such that $\displaystyle\lim_{n \to \infty} u_n = 0$ and $\displaystyle\lim_{n \to \infty} \frac{u_n - u_{n+1}}{u_n^{\alpha}}$ exists and is nonzero. Prove that $\displaystyle\sum_{n=1}^{\infty} u_n$ converges if and only if $\alpha < 2$.
\end{problembox}

\stepcounter{probnum}
\begin{problembox}
\textbf{\theprobnum.} Prove that, for any two bounded functions $g_1, g_2 : \mathbb{R} \to [1,\infty)$, there exist functions $h_1, h_2 : \mathbb{R} \to \mathbb{R}$ such that for every $x \in \mathbb{R}$,
\[
\sup_{s \in \mathbb{R}} \left(g_1(s)^x g_2(s)\right) = \max_{t \in \mathbb{R}} \left(x h_1(t) + h_2(t)\right).
\]
\end{problembox}

\stepcounter{probnum}
\begin{problembox}
\textbf{\theprobnum.} Is there a strictly increasing function $f : \mathbb{R} \to \mathbb{R}$ such that $f'(x) = f(f(x))$ for all reals $x$?
\end{problembox}

\stepcounter{probnum}
\begin{problembox}
\textbf{\theprobnum.} Does there exist a continuously differentiable function $f : \mathbb{R} \to \mathbb{R}$ satisfying $f(x) > 0$ and $f'(x) = f(f(x))$ for all reals $x$?
\end{problembox}

\stepcounter{probnum}
\begin{problembox}
\textbf{\theprobnum.} Find all differentiable functions $f : \mathbb{R} \to \mathbb{R}$ such that
\[
f'(x) = \frac{f(x+n) - f(x)}{n}
\]
for all real numbers $x$ and all positive integers $n$.
\end{problembox}

\stepcounter{probnum}
\begin{problembox}
\textbf{\theprobnum.} Find the set of all real numbers $k$ with the following property: for any positive, differentiable function $f$ that satisfies $f'(x) > f(x)$ for all $x$, there is some number $N$ such that $f(x) > e^{kx}$ for all $x > N$.
\end{problembox}

\stepcounter{probnum}
\begin{problembox}
\textbf{\theprobnum.} Let $f : (1,\infty) \to \mathbb{R}$ be a differentiable function such that
\[
f'(x) = \frac{x^2 - (f(x))^2}{x^2\left((f(x))^2+1\right)} \qquad \text{for all } x > 1.
\]
Prove that $\displaystyle\lim_{x \to \infty} f(x) = \infty$.
\end{problembox}

\stepcounter{probnum}
\begin{problembox}
\textbf{\theprobnum.} Define $f : \mathbb{R} \to \mathbb{R}$ by
\[
f(x) = \begin{cases} x & \text{if } x \le e \\ x f(\ln x) & \text{if } x > e \end{cases}
\]
Does $\displaystyle\sum_{n=1}^{\infty} \frac{1}{f(n)}$ converge?
\end{problembox}

\stepcounter{probnum}
\begin{problembox}
\textbf{\theprobnum.} Find all sequences $a_0, a_1, \ldots, a_n$ of real numbers with $a_n \ne 0$, for which the following statement is true: if $f : \mathbb{R} \to \mathbb{R}$ is an $n$ times differentiable function and $x_0 < x_1 < \cdots < x_n$ are real numbers such that $f(x_0) = f(x_1) = \cdots = f(x_n) = 0$, then there is $h \in (x_0,x_n)$ for which
\[
a_0 f(h) + a_1 f'(h) + \cdots + a_n f^{(n)}(h) = 0.
\]
\end{problembox}

\stepcounter{probnum}
\begin{problembox}
\textbf{\theprobnum.} Let $f : (0,\infty) \to \mathbb{R}$ be a differentiable function. Assume that
\[
\lim_{x \to \infty} \left(f(x) + \frac{f'(x)}{x}\right) = 0.
\]
Prove that $\displaystyle\lim_{x \to \infty} f(x) = 0$.
\end{problembox}

\stepcounter{probnum}
\begin{problembox}
\textbf{\theprobnum.} Let $f$ be a twice continuously differentiable function on $(0,\infty)$ such that $\displaystyle\lim_{x \to 0^+} f'(x) = -\infty$ and $\displaystyle\lim_{x \to 0^+} f''(x) = +\infty$. Show that
\[
\lim_{x \to 0^+} \frac{f(x)}{f'(x)} = 0.
\]
\end{problembox}

\stepcounter{probnum}
\begin{problembox}
\textbf{\theprobnum.} Determine all Riemann integrable functions $f : \mathbb{R} \to \mathbb{R}$ such that
\[
\int_0^{x+1/n} f(t)\,dt = \int_0^x f(t)\,dt + \frac{1}{n} f(x),
\]
for all reals $x$ and all positive integers $n$.
\end{problembox}

\stepcounter{probnum}
\begin{problembox}
\textbf{\theprobnum.} For each continuous function $f : [0,1] \to \mathbb{R}$, let $I(f) = \displaystyle\int_0^1 x^2 f(x)\,dx$ and $J(f) = \displaystyle\int_0^1 x f^2(x)\,dx$. Find the maximum value of $I(f) - J(f)$ over all such functions.
\end{problembox}

\stepcounter{probnum}
\begin{problembox}
\textbf{\theprobnum.} Let $f : [0,1] \to [0,\infty)$ be a continuous function. Show that
\[
\int_0^1 f^3(x)\,dx \ge 4 \int_0^1 x^2 f(x)\,dx \int_0^1 x f^2(x)\,dx.
\]
\end{problembox}

\stepcounter{probnum}
\begin{problembox}
\textbf{\theprobnum.} Find all differentiable functions $f : [0,1] \to \mathbb{R}$ with continuous derivative such that $f(0) = -\dfrac{1}{6}$ and
\[
\int_0^1 (f'(x))^2\,dx \le 2 \int_0^1 f(x)\,dx.
\]
\end{problembox}

\stepcounter{probnum}
\begin{problembox}
\textbf{\theprobnum.} a) Let $f : [0,1] \to \mathbb{R}$ be continuously differentiable and satisfy $f(1) = 0$. Show that
\[
\int_0^1 |f(x)|^2\,dx \le 4 \int_0^1 x^2 (f'(x))^2\,dx.
\]
b) Let $f : [0,1] \to \mathbb{R}$ be a differentiable function with continuous derivative such that $f(1) = 0$. Show that
\[
4 \int_0^1 x^2 |f'(x)|^2\,dx \ge \int_0^1 |f(x)|^2\,dx + \left(\int_0^1 f(x)\,dx\right)^2.
\]
\end{problembox}

\stepcounter{probnum}
\begin{problembox}
\textbf{\theprobnum.} Let $f : [0,1] \to [0,1)$ be a continuous function such that
\[
\int_0^1 f(x^n)\,dx \le \int_0^1 f^n(x)\,dx, \qquad n \ge 1.
\]
Show that $f(x) = 0$ for all $x \in [0,1]$.
\end{problembox}

\stepcounter{probnum}
\begin{problembox}
\textbf{\theprobnum.} Find all continuous functions $f : \mathbb{R} \to [1,\infty)$ for which there exists $a \in \mathbb{R}$ and a positive integer $k$ such that
\[
f(x) f(2x) \cdots f(nx) \le a n^k,
\]
for every real number $x$ and positive integer $n$.
\end{problembox}

\stepcounter{probnum}
\begin{problembox}
\textbf{\theprobnum.} Let $f : [0,1] \to \mathbb{R}^+$ be an integrable function. Compute the following limits:\\
a) $\displaystyle\lim_{n \to \infty} \prod_{k=1}^{n} \left(1+\frac{1}{n} f\!\left(\frac{k}{n}\right)\right)$;\\
b) $\displaystyle\lim_{n \to \infty} \left(\sum_{k=1}^{n} \exp\!\left(\frac{1}{n} f\!\left(\frac{k}{n}\right)\right) - n\right)$.
\end{problembox}

\stepcounter{probnum}
\begin{problembox}
\textbf{\theprobnum.} Prove that
\[
\lim_{n \to \infty} n \left(\frac{\pi}{4} - n \int_0^1 \frac{x^n}{1+x^{2n}}\,dx\right) = \int_0^1 f(x)\,dx,
\]
where $f(x) = \dfrac{\arctan x}{x}$ if $x \in (0,1]$ and $f(0) = 1$.
\end{problembox}

\stepcounter{probnum}
\begin{problembox}
\textbf{\theprobnum.} Let $f : [0,1] \to \mathbb{R}$ be a continuous function such that $\displaystyle\int_0^1 f(x)\,dx = 0$. Prove that there is $c \in (0,1)$ such that
\[
\int_0^c x f(x)\,dx = 0.
\]
\end{problembox}

\stepcounter{probnum}
\begin{problembox}
\textbf{\theprobnum.} Let $f : [-1,1] \to \mathbb{R}$ be a continuous function. Show that:\\
a) if $\displaystyle\int_0^1 f(\sin(x+\alpha))\,dx = 0$ for every $\alpha \in \mathbb{R}$, then $f(x) = 0$ for all $x \in [-1,1]$;\\
b) if $\displaystyle\int_0^1 f(\sin(nx))\,dx = 0$ for every $n \in \mathbb{Z}$, then $f(x) = 0$ for all $x \in [-1,1]$.
\end{problembox}

\stepcounter{probnum}
\begin{problembox}
\textbf{\theprobnum.} Let $f : [0,\infty) \to \mathbb{R}$ be a periodic function, with period $1$, integrable on $[0,1]$. For a strictly increasing and unbounded sequence $(x_n)_{n \ge 0}$, $x_0 = 0$, with $\displaystyle\lim_{n \to \infty} (x_{n+1}-x_n) = 0$, denote $r(n) = \max\{k \mid x_k \le n\}$.\\
a) Show that
\[
\lim_{n \to \infty} \frac{1}{n} \sum_{k=1}^{r(n)} (x_k - x_{k+1}) f(x_k) = \int_0^1 f(x)\,dx.
\]
b) Show that
\[
\lim_{n \to \infty} \frac{1}{\log n} \sum_{k=1}^{r(n)} \frac{f(\log k)}{k} = \int_0^1 f(x)\,dx.
\]
\end{problembox}

\stepcounter{probnum}
\begin{problembox}
\textbf{\theprobnum.} Let $f$ and $g$ be two continuous, distinct functions from $[0,1]$ to $(0,+\infty)$ such that $\displaystyle\int_0^1 f(x)\,dx = \int_0^1 g(x)\,dx$. Let $(y_n)_{n \ge 1}$ be the sequence defined by
\[
y_n = \int_0^1 \frac{f^{n+1}(x)}{g^n(x)}\,dx.
\]
Prove that $(y_n)$ is an increasing and divergent sequence.
\end{problembox}

\stepcounter{probnum}
\begin{problembox}
\textbf{\theprobnum.} Let $f : \mathbb{R} \to \mathbb{R}$ be a continuous and bounded function such that
\[
x \int_x^{x+1} f(t)\,dt = \int_0^x f(t)\,dt, \qquad \text{for any } x \in \mathbb{R}.
\]
Prove that $f$ is a constant function.
\end{problembox}

\stepcounter{probnum}
\begin{problembox}
\textbf{\theprobnum.} Let $f : [0,1] \to \mathbb{R}$ be an integrable function such that $0 < \left|\displaystyle\int_0^1 f(x)\,dx\right| \le 1$. Show that there exist $x_1 \ne x_2$ in $[0,1]$ such that
\[
\int_{x_1}^{x_2} f(x)\,dx = (x_1-x_2)^{2002}.
\]
\end{problembox}

\stepcounter{probnum}
\begin{problembox}
\textbf{\theprobnum.} Let $f : \mathbb{R} \to (0,\infty)$ be a continuously differentiable function. Prove that
\[
\left|\int_0^1 f^3(x)\,dx - f^2(0) \int_0^1 f(x)\,dx\right| \le \max_{[0,1]} |f'(x)| \left(\int_0^1 f(x)\,dx\right)^2.
\]
\end{problembox}

\stepcounter{probnum}
\begin{problembox}
\textbf{\theprobnum.} Let $f : [0,1] \to [0,\infty)$ be an integrable function. Show that
\[
\int_0^1 f(x)\,dx \int_0^1 f^2(x)\,dx \le \int_0^1 f^3(x)\,dx,
\]
and
\[
\int_0^1 x f(x)\,dx \int_0^1 x^2 f(x)\,dx \le \int_0^1 f(x)\,dx \int_0^1 x^3 f(x)\,dx.
\]
\end{problembox}

\stepcounter{probnum}
\begin{problembox}
\textbf{\theprobnum.} Show that for any continuous function $f : [0,1] \to \mathbb{R}$,
\[
\frac{1}{4} \int_0^1 f^2(x)\,dx + 2 \left(\int_0^1 f(x)\,dx\right)^2 \ge 3 \int_0^1 f(x)\,dx \int_0^1 x f(x)\,dx.
\]
\end{problembox}

\stepcounter{probnum}
\begin{problembox}
\textbf{\theprobnum.} Let $f : [0,1] \to \mathbb{R}$ be an integrable function such that
\[
\int_0^1 f(x)\,dx = \int_0^1 x f(x)\,dx = \int_0^1 x^2 f(x)\,dx = 1.
\]
Show that $\displaystyle\int_0^1 f^2(x)\,dx \ge 9$.
\end{problembox}

\stepcounter{probnum}
\begin{problembox}
\textbf{\theprobnum.} Let $f : [0,1] \to \mathbb{R}$ be a differentiable function with continuous derivative such that
\[
\int_0^1 f(x)\,dx = \int_0^1 x f(x)\,dx = 1.
\]
Show that $\displaystyle\int_0^1 (f'(x))^2\,dx \ge 30$.
\end{problembox}

\stepcounter{probnum}
\begin{problembox}
\textbf{\theprobnum.} a) Consider $f : [0,\infty) \to \mathbb{R}$ and $g : [0,1] \to \mathbb{R}$ two continuous functions such that $\displaystyle\lim_{x \to \infty} f(x) = L \in \mathbb{R}$. Show that
\[
\lim_{n \to \infty} \frac{1}{n} \int_0^n f(x) g\!\left(\frac{x}{n}\right) dx = L \int_0^1 g(x)\,dx.
\]
b) Find $\displaystyle\lim_{n \to \infty} \frac{1}{n} \int_0^n \frac{x \log(1+x/n)}{1+x}\,dx$.
\end{problembox}

\stepcounter{probnum}
\begin{problembox}
\textbf{\theprobnum.} a) Let $f : [0,1] \to \mathbb{R}$ be a continuous function such that $\displaystyle\int_0^1 f(x)\,dx = 0$. Show that there exists $c \in (0,1)$ such that
\[
f(c) = \int_0^c f(x)\,dx.
\]
b) Let $f : [0,1] \to \mathbb{R}$ be a continuous function such that $f(1) = 0$. Show that there exists $c \in (0,1)$ such that
\[
f(c) = \int_0^c f(x)\,dx.
\]
\end{problembox}

\stepcounter{probnum}
\begin{problembox}
\textbf{\theprobnum.} Let $f : [0,1] \to \mathbb{R}$ be a differentiable function for which there exist $a, b \in (0,1)$, $a < b$, such that $\displaystyle\int_0^a f(x)\,dx = 0$ and $\displaystyle\int_b^1 f(x)\,dx = 0$. Let $M = \displaystyle\sup_{x \in [0,1]} |f'(x)|$. Show that
\[
\left|\int_0^1 f(x)\,dx\right| \le \frac{1-a+b}{4} M.
\]
Moreover, if $f : [0,1] \to [0,\infty)$ is differentiable on $[0,1]$ and there exists $a \in (0,1]$ with $\displaystyle\int_0^a f(x)\,dx = 0$, show that
\[
\left|\int_0^1 f(x)\,dx\right| \le \frac{1-a}{2} \sup_{x \in (0,1)} |f'(x)|.
\]
\end{problembox}

\stepcounter{probnum}
\begin{problembox}
\textbf{\theprobnum.} a) Let $f_0 : [0,1] \to \mathbb{R}$ be a continuous function. Define the sequence of functions $f_n : [0,1] \to \mathbb{R}$ by $f_n(x) = \displaystyle\int_0^x f_{n-1}(t)\,dt$ for all integers $n \ge 1$. Prove that the series $\displaystyle\sum_{n=1}^{\infty} f_n(x)$ is convergent for every $x \in [0,1]$ and find an explicit formula for the sum of the series.\\
b) Let $F_0 = \ln x$. For $n \ge 0$ and $x > 0$, let $F_{n+1}(x) = \displaystyle\int_0^x F_n(t)\,dt$. Evaluate $\displaystyle\lim_{n \to \infty} \frac{n! F_n(1)}{\ln n}$.
\end{problembox}

\stepcounter{probnum}
\begin{problembox}
\textbf{\theprobnum.} Let $f : [0,\infty) \to \mathbb{R}$ be a strictly decreasing continuous function such that $\displaystyle\lim_{x \to \infty} f(x) = 0$. Prove that $\displaystyle\int_0^{\infty} \frac{f(x)-f(x+1)}{f(x)}\,dx$ diverges.
\end{problembox}

\stepcounter{probnum}
\begin{problembox}
\textbf{\theprobnum.} Let $(a_n)_{n \ge 1}$ be defined by
\[
a_n = \frac{3}{2} \log n - \int_0^1 \log\left(1^t + 2^t + \cdots + n^t\right) dt.
\]
i) Show that the sequence $a_n$ is convergent and find its limit $a$.\\
ii) Show that $0 < n(a-a_n) < \dfrac{3}{2}$.
\end{problembox}

\stepcounter{probnum}
\begin{problembox}
\textbf{\theprobnum.} Let $f : \mathbb{R} \to \mathbb{R}$ be a continuous and periodic function of period $T > 0$. Show that:\\
i) $\displaystyle\lim_{x \to \infty} \frac{1}{x} \int_0^x f(t)\,dt = \frac{1}{T} \int_0^T f(t)\,dt$.\\
ii) $\displaystyle\lim_{n \to \infty} \int_a^b f(nx)\,dx = \frac{b-a}{T} \int_a^b f(t)\,dt$.\\
iii) Let $f, g : \mathbb{R} \to \mathbb{R}$ be continuous functions such that $f(x+1)=f(x)$ and $g(x+1)=g(x)$ for all real $x$. Prove that
\[
\lim_{n \to \infty} \int_0^1 f(x) g(nx)\,dx = \int_0^1 f(x)\,dx \int_0^1 g(x)\,dx.
\]
\end{problembox}

\stepcounter{probnum}
\begin{problembox}
\textbf{\theprobnum.} Let $a_1, a_2, \ldots$ be real numbers. Suppose there is a constant $A$ such that for all $n$,
\[
\int_{-\infty}^{\infty} \left(\sum_{i=1}^{n} \frac{1}{1+(x-a_i)^2}\right)^2 dx \le An.
\]
Prove there is a constant $B > 0$ such that for all $n$,
\[
\sum_{i,j=1}^{n} \left(1+(a_i-a_j)^2\right) \ge Bn^3.
\]
\end{problembox}

\stepcounter{probnum}
\begin{problembox}
\textbf{\theprobnum.} a) (Hermite--Hadamard) Let $f : [a,b] \to \mathbb{R}$ be a convex function. Show that
\[
(b-a) f\left(\frac{a+b}{2}\right) \le \int_a^b f(x)\,dx \le (b-a) \frac{f(a)+f(b)}{2}.
\]
b) Show that $\displaystyle\int_k^{k+1} \log x\,dx \ge \frac{\log k + \log(k+1)}{2}$ for $k \ge 1$, and $\displaystyle\int_{k-1/2}^{k+1/2} \log x\,dx \le \log k$ for $k \ge 1$.\\
c) Consider the sequence $(a_n)_{n \ge 1}$ defined by
\[
a_n = \int_1^n \log x\,dx - \log 2 - \cdots - \log(n-1) - \frac{1}{2}\log n, \qquad n \ge 1.
\]
Show that $a_n$ is increasing and $0 \le a_n \le \dfrac{1}{2}\log \dfrac{5}{4}$.\\
d) Prove that $\sqrt[4]{\dfrac{4}{5}}\, e\sqrt{n}\left(\dfrac{n}{e}\right)^n \le n! \le e\sqrt{n}\left(\dfrac{n}{e}\right)^n$ for $n \ge 1$.\\
e) (Stirling) Show that $\displaystyle\lim_{n \to \infty} \dfrac{n!}{\left(\frac{n}{e}\right)^n \sqrt{2\pi n}} = 1$.
\end{problembox}

\stepcounter{probnum}
\begin{problembox}
\textbf{\theprobnum.} Show that for any continuous function $f : [-1,1] \to \mathbb{R}$ we have the inequalities
\[
\int_{-1}^{1} f^2(x)\,dx \ge \frac{1}{2}\left(\int_{-1}^{1} f(x)\,dx\right)^2 + \frac{3}{2}\left(\int_{-1}^{1} x f(x)\,dx\right)^2,
\]
and
\[
\int_{-1}^{1} f^2(x)\,dx \ge \frac{5}{2}\left(\int_{-1}^{1} x^2 f(x)\,dx\right)^2 + \frac{3}{2}\left(\int_{-1}^{1} x f(x)\,dx\right)^2.
\]
\end{problembox}

\stepcounter{probnum}
\begin{problembox}
\textbf{\theprobnum.} a) Let $f : [0,1] \to \mathbb{R}$ be a continuous function. Show that $\displaystyle\lim_{n \to \infty} \int_0^1 x^n f(x)\,dx = 0$.\\
b) Let $f : [0,1] \to \mathbb{R}$ be a continuous function. Show that
\[
\lim_{n \to \infty} n \int_0^1 x^n f(x)\,dx = f(1), \qquad \lim_{n \to \infty} n \int_0^1 x^n f(x^n)\,dx = \int_0^1 f(x)\,dx.
\]
c) Let $f, g : [0,1] \to \mathbb{R}$ be two continuous functions. Show that
\[
\lim_{n \to \infty} \frac{\int_0^1 x^n f(x)\,dx}{\int_0^1 x^n g(x)\,dx} = \frac{f(1)}{g(1)}, \qquad \lim_{n \to \infty} \frac{\int_0^1 x^n f(x^n)\,dx}{\int_0^1 x^n g(x^n)\,dx} = \frac{\int_0^1 f(x)\,dx}{\int_0^1 g(x)\,dx}.
\]
d) Find a real number $c$ and a positive number $L$ for which
\[
\lim_{r \to \infty} r^c \frac{\int_0^{\pi/2} x^r \sin x\,dx}{\int_0^{\pi/2} x^r \cos x\,dx} = L.
\]
\end{problembox}

\stepcounter{probnum}
\begin{problembox}
\textbf{\theprobnum.} Let $(a_n)_{n \in \mathbb{N}}$ be the sequence defined by
\[
a_0 = 1, \qquad a_{n+1} = \frac{1}{n+1} \sum_{k=0}^{n} \frac{a_k}{n-k+2}.
\]
Find the limit $\displaystyle\lim_{n \to \infty} \sum_{k=0}^{n} \frac{a_k}{2^k}$.
\end{problembox}

\stepcounter{probnum}
\begin{problembox}
\textbf{\theprobnum.} Prove that for any real numbers $a_1, a_2, \ldots, a_n$,
\[
\sum_{1 \le i,j \le n} \frac{ij}{i+j-1} a_i a_j \ge \left(\sum_{i=1}^{n} a_i\right)^2.
\]
\end{problembox}

\stepcounter{probnum}
\begin{problembox}
\textbf{\theprobnum.} (Hardy) Let $p > 1$ and let $f : [0,\infty) \to \mathbb{R}$ be a differentiable increasing function such that $f(0) = 0$ and $\displaystyle\int_0^{\infty} (f'(x))^p\,dx$ is finite. Show that
\[
\int_0^{\infty} x^{-p} |f(x)|^p\,dx \le \left(\frac{p}{p-1}\right)^p \int_0^{\infty} (f'(x))^p\,dx.
\]
\end{problembox}

\stepcounter{probnum}
\begin{problembox}
\textbf{\theprobnum.} Let $u_1, u_2, \ldots, u_n \in C([0,1]^n)$ be nonnegative and continuous functions, where $u_j$ does not depend on the $j$-th variable for $j = 1, 2, \ldots, n$. Show that
\[
\left(\int_{[0,1]^n} \prod_{j=1}^{n} u_j\right)^{n-1} \le \prod_{j=1}^{n} \int_{[0,1]^n} u_j^{n-1}.
\]
\end{problembox}

\stepcounter{probnum}
\begin{problembox}
\textbf{\theprobnum.} (Zagier) Prove that if $a_1, a_2, \ldots, a_n, b_1, b_2, \ldots, b_n$ are nonnegative, then
\[
\left(\sum_{1 \le i,j \le n} \min(a_i,a_j)\right) \left(\sum_{1 \le i,j \le n} \min(b_i,b_j)\right) \ge \left(\sum_{1 \le i,j \le n} \min(a_i,b_j)\right)^2.
\]
\end{problembox}

\stepcounter{probnum}
\begin{problembox}
\textbf{\theprobnum.} Let $\{D_1, D_2, \ldots, D_n\}$ be a set of disks in the Euclidean plane and $a_{ij} = S(D_i \cap D_j)$ be the area of $D_i \cap D_j$. Prove that for any real numbers $x_1, x_2, \ldots, x_n$,
\[
\sum_{i=1}^{n} \sum_{j=1}^{n} a_{ij} x_i x_j \ge 0.
\]
\end{problembox}

\stepcounter{probnum}
\begin{problembox}
\textbf{\theprobnum.} Let $a, b \in \mathbb{R}$, $f : [a,b] \to [0,1]$ an increasing function, and define the sequence $(c_n)_{n \ge 1}$ by $c_n = n \displaystyle\int_a^b f^n(x)\,dx$. Show that:\\
i) $\displaystyle\lim_{n \to \infty} (c_{n+1}-c_n) = 0$.\\
ii) $\displaystyle\lim_{n \to \infty} \frac{c_{n+1}}{c_n} = f(b^-)$.
\end{problembox}

\stepcounter{probnum}
\begin{problembox}
\textbf{\theprobnum.} Let $f : [0,1] \to \mathbb{R}$ be a $C^1$ function. Show that
\[
\lim_{n \to \infty} n \left[\frac{1}{2}\left(\int_0^1 f(x)\,dx\right)^2 - \frac{1}{n^2} \sum_{1 \le i,j \le n} f\!\left(\frac{i}{n}\right) f\!\left(\frac{j}{n}\right)\right] = \frac{1}{2}\int_0^1 f^2(x)\,dx + \frac{f(0)-f(1)}{2}\int_0^1 f(x)\,dx.
\]
\end{problembox}

\stepcounter{probnum}
\begin{problembox}
\textbf{\theprobnum.} a) Find the limit $\displaystyle\lim_{n \to \infty} \left(\int_0^1 (1+x^n)^n\,dx\right)^{1/n}$.\\
b) Prove that $\displaystyle\lim_{n \to \infty} n^2 \left(\int_0^1 \sqrt[n]{1+x^n}\,dx - 1\right) = \dfrac{\pi^2}{12}$.
\end{problembox}

\stepcounter{probnum}
\begin{problembox}
\textbf{\theprobnum.} Let $f : [0,1] \to \mathbb{R}$ be a differentiable function such that $f'$ is bounded. Show that
\[
\int_0^1 f^2(x)\,dx - \left(\int_0^1 f(x)\,dx\right)^2 \le \frac{1}{12} \left(\sup_{x \in [0,1]} |f'(x)|\right)^2.
\]
\end{problembox}

\stepcounter{probnum}
\begin{problembox}
\textbf{\theprobnum.} Assume that $f : \mathbb{R} \to \mathbb{R}$ is a continuous injection. Prove that if there exists $n$ such that the $n$-th iterate of $f$ is the identity, that is, $f^n(x) = x$ for all $x \in \mathbb{R}$, then:\\
a) $f(x) = x$ for all $x \in \mathbb{R}$, if $f$ is strictly increasing;\\
b) $f^2(x) = x$ for all $x \in \mathbb{R}$, if $f$ is strictly decreasing.
\end{problembox}

\stepcounter{probnum}
\begin{problembox}
\textbf{\theprobnum.} (Continuity of square-root-type functional equations)\\
a) Assume that $f : \mathbb{R} \to \mathbb{R}$ satisfies $f(x)^2 = -x$ for all $x \in \mathbb{R}$. Show that $f$ cannot be continuous.\\
b) Find all continuous functions $f : \mathbb{R} \to \mathbb{R}$ such that $f^2(x) = \cos(x)$ for all $x \in \mathbb{R}$.
\end{problembox}

\stepcounter{probnum}
\begin{problembox}
\textbf{\theprobnum.} Find all functions $f : \mathbb{R} \to \mathbb{R}$ which have the Intermediate Value Property and such that there is $n \in \mathbb{N}$ for which $f^n(x) = -x$ for all $x \in \mathbb{R}$.
\end{problembox}

\stepcounter{probnum}
\begin{problembox}
\textbf{\theprobnum.} Give an example of a continuous function $f : \mathbb{R} \to \mathbb{R}$ which attains each of its values exactly three times. Does there exist a continuous function $f : \mathbb{R} \to \mathbb{R}$ which attains each of its values exactly two times?
\end{problembox}

\stepcounter{probnum}
\begin{problembox}
\textbf{\theprobnum.} Let $f \in C^2[0,1]$ be such that $f(0) = f(1)$. Show that
\[
3 f'(1)^2 \le \int_0^1 f''(x)^2\,dx.
\]
\end{problembox}

\stepcounter{probnum}
\begin{problembox}
\textbf{\theprobnum.} (Monotonic functions and continuity)\\
a) If $f$ is monotonic on $[a,b]$, then the set of discontinuities of $f$ is countable.\\
b) Let $f : \mathbb{R} \to \mathbb{R}$ be a surjective monotonous function. Show that $f$ must be continuous.
\end{problembox}

\stepcounter{probnum}
\begin{problembox}
\textbf{\theprobnum.} Let $f : [0,1] \to [0,1]$ be a continuous function. Consider any sequence given by $x_{n+1} = f(x_n)$. Show that $(x_n)$ converges if and only if $\lim_{n \to \infty} (x_{n+1} - x_n) = 0$.
\end{problembox}

\stepcounter{probnum}
\begin{problembox}
\textbf{\theprobnum.} Let $f$ be a three times continuously differentiable function with the property that $f(0) = f'(0) = 0 < f''(0)$. Define
\[
g(x) = \frac{d}{dx} \left( \frac{\sqrt{f(x)}}{f'(x)} \right)
\]
for $x \ne 0$ and $g(0) = 0$. Show that $g$ is bounded in a neighbourhood of $0$. Does this conclusion still follow if $f$ is only twice continuously differentiable?
\end{problembox}

\stepcounter{probnum}
\begin{problembox}
\textbf{\theprobnum.} Let $f : [0,1] \to \mathbb{R}$ be given by $f(x) = 2x(1-x)$. Denote $f_n = f \circ \cdots \circ f$ for the $n$-fold composition of $f$. Determine
\[
\int_0^1 f_n(x)\,dx.
\]
Does this sequence converge?
\end{problembox}

\stepcounter{probnum}
\begin{problembox}
\textbf{\theprobnum.} Let $f : \mathbb{R} \to \mathbb{R}$ be a twice differentiable function satisfying $f(0) = 2$, $f'(0) = -2$ and $f(1) = 1$. Prove that there exists some $c \in (0,1)$ such that
\[
f(c) \cdot f'(c) + f''(c) = 0.
\]
\end{problembox}

\stepcounter{probnum}
\begin{problembox}
\textbf{\theprobnum.} Let $c \in (0,1)$ and define
\[
f(x) =
\begin{cases}
x/c & \text{if } x \in [0,c]; \\
(1-x)/(1-c) & \text{if } x \in [c,1].
\end{cases}
\]
Show that the equation $f^n(x) = x$ only has finitely many solutions for every $n \ge 1$, where $f^n$ denotes the $n$-fold composition of $f$ with itself. Show also that for every $n$, there are solutions to $f^n(x) = x$ for which this $n$ is minimal.
\end{problembox}

\stepcounter{probnum}
\begin{problembox}
\textbf{\theprobnum.} (*) Let $f : (0,1) \to [0,\infty)$ be a function that is non-zero only on the distinct points $a_1, a_2, \ldots$. Suppose that
\[
\sum_{n=1}^{\infty} f(a_n) < \infty.
\]
Prove that $f$ is differentiable at at least one $x \in (0,1)$.
\end{problembox}

\stepcounter{probnum}
\begin{problembox}
\textbf{\theprobnum.} Let $f : [0,1] \to [0,1]$ be a strictly increasing function. Show that for some $x \in [0,1]$, it holds that $f(x) = x$. What if $f$ is strictly decreasing?
\end{problembox}

\stepcounter{probnum}
\begin{problembox}
\textbf{\theprobnum.} Let $f : [a,b] \to [a,b]$ be a continuous function and let $p \in [a,b]$. Define the sequence $(p_n)$ by $p_0 = p$ and $p_{n+1} = f(p_n)$ for all $n \ge 0$. Suppose that the set $T_p = \{p_n \mid n \in \mathbb{N}\}$ is closed. Prove that $T_p$ is finite.
\end{problembox}

\stepcounter{probnum}
\begin{problembox}
\textbf{\theprobnum.} (*) Let $g : [0,1] \to \mathbb{R}$ be a function and define the sequence of functions $f_n : (0,1] \to \mathbb{R}$ by setting $f_0(x) = g(x)$ and
\[
f_{n+1}(x) = \frac{1}{x} \int_0^x f_n(t)\,dt
\]
for all $x \in (0,1]$ and $n \ge 0$. Determine $\lim_{n \to \infty} f_n(x)$ for every $x \in (0,1]$.
\end{problembox}

\stepcounter{probnum}
\begin{problembox}
\textbf{\theprobnum.} Let $f, g : [a,b] \to [0,\infty)$ be continuous and non-decreasing functions such that for each $x \in [a,b]$, we have
\[
\int_a^x \sqrt{f(t)}\,dt \le \int_a^x \sqrt{g(t)}\,dt,
\]
with equality for $x = b$. Prove that
\[
\int_a^b \sqrt{1+f(t)}\,dt \ge \int_a^b \sqrt{1+g(t)}\,dt.
\]
\end{problembox}

\stepcounter{probnum}
\begin{problembox}
\textbf{\theprobnum.} Let $f(x) = x^2 + bx + c$ where $b, c \in \mathbb{R}$. Now define the possibly empty set $M = \{x \in \mathbb{R} : |f(x)| < 1\}$. Clearly $M$ is the union of disjoint open intervals. Let $|M|$ denote the sum of their lengths. Prove that $|M| \le 2\sqrt{2}$.
\end{problembox}

\stepcounter{probnum}
\begin{problembox}
\textbf{\theprobnum.} Let $f : \mathbb{R} \to \mathbb{R}$ be a three times differentiable function. Prove that for some $c \in (-1,1)$, it holds that
\[
f'''(c) = 3f(1) - 3f(-1) - 6f'(0).
\]
\end{problembox}

\stepcounter{probnum}
\begin{problembox}
\textbf{\theprobnum.} Determine all functions $f : \mathbb{R} \to \mathbb{R}$ such that for any real numbers $a < b$, the image $f([a,b])$ is a closed interval of length $b - a$.
\end{problembox}

\stepcounter{probnum}
\begin{problembox}
\textbf{\theprobnum.} Let $f : \mathbb{R} \to \mathbb{R}$ be a continuous function with the property that for every $c > 0$, the graph of $f$ can be transformed into the graph of $c \cdot f$ by translating and rotating. Does this imply that $f(x) = ax + b$ for some $a, b \in \mathbb{R}$?
\end{problembox}

\stepcounter{probnum}
\begin{problembox}
\textbf{\theprobnum.} Let $C \subset \mathbb{R}$ be a compact set and let $f : C \to C$ be a non-decreasing continuous function. Show that for some $p \in C$ it holds that $f(p) = p$.
\end{problembox}

\stepcounter{probnum}
\begin{problembox}
\textbf{\theprobnum.} Determine all continuous functions $f : \mathbb{R} \to \mathbb{R}$ with the property that
\[
x - y \in \mathbb{Q} \implies f(x) - f(y) \in \mathbb{Q}.
\]
\end{problembox}

\stepcounter{probnum}
\begin{problembox}
\textbf{\theprobnum.} Let $f, g : \mathbb{R} \to \mathbb{R}$ be two functions with the property that $f(r) \le g(r)$ for all $r \in \mathbb{Q}$. If $f$ and $g$ are non-decreasing, does this imply that $f(x) \le g(x)$ for all $x \in \mathbb{R}$? What if $f$ and $g$ are both continuous?
\end{problembox}

\stepcounter{probnum}
\begin{problembox}
\textbf{\theprobnum.} Let $f : \mathbb{R} \to \mathbb{R}$ be a continuous function. We say some $x \in \mathbb{R}$ is a shadow-point of $f$ if for some $y > x$ it holds that $f(y) > f(x)$. Let $a < b$ now be two real numbers that are not shadow-points of $f$, with the property that all points in $(a,b)$ are shadow-points of $f$. Show that $f(x) \le f(b)$ for all $x \in (a,b)$ and that $f(a) = f(b)$.
\end{problembox}

\stepcounter{probnum}
\begin{problembox}
\textbf{\theprobnum.} Compute
\[
\lim_{A \to \infty} \frac{1}{A} \int_1^A A^{1/x}\,dx.
\]
\end{problembox}

\stepcounter{probnum}
\begin{problembox}
\textbf{\theprobnum.} Let $f : [a,b] \to \mathbb{R}$ be a continuous function that is differentiable on $(a,b)$. Suppose that $f$ has infinitely many zeroes, but that there is no $x \in (a,b)$ such that $f(x) = f'(x) = 0$. Prove that $f(a) f(b) = 0$ and give an example of such a function on $[0,1]$.
\end{problembox}

\stepcounter{probnum}
\begin{problembox}
\textbf{\theprobnum.} Let $f : [0,\infty) \to \mathbb{R}$ be a continuous function satisfying the property that $\lim_{x \to \infty} f(x) = L$ exists. Prove that
\[
\lim_{n \to \infty} \int_0^1 f(nx)\,dx = L.
\]
\end{problembox}

\stepcounter{probnum}
\begin{problembox}
\textbf{\theprobnum.} (*) Define the sequence $f_1, f_2, \ldots : [0,1] \to \mathbb{R}$ of continuously differentiable functions by
\[
f_1 = 1, \qquad f_{n+1}' = f_n f_{n+1} \text{ on } (0,1) \qquad \text{and} \qquad f_{n+1}(0) = 1.
\]
Show that $\lim_{n \to \infty} f_n(x)$ exists for every $x \in [0,1]$ and determine the limit function.
\end{problembox}

\stepcounter{probnum}
\begin{problembox}
\textbf{\theprobnum.} Let $f, g : \mathbb{R} \to \mathbb{R}$ be continuous functions such that $g$ is differentiable. Suppose that
\[
(f(0) - g'(0))(g'(1) - f(1)) > 0.
\]
Show that there exists some $c \in (0,1)$ such that $f(c) = g'(c)$.
\end{problembox}

\stepcounter{probnum}
\begin{problembox}
\textbf{\theprobnum.} Let $V$ be the set of all continuous functions $f : [0,1] \to \mathbb{R}$, differentiable on $(0,1)$, with the property that $f(0) = 0$ and $f(1) = 1$. Determine all $\alpha \in \mathbb{R}$ such that for every $f \in V$, there exists some $\xi \in (0,1)$ such that $f(\xi) + \alpha = f'(\xi)$.
\end{problembox}

\stepcounter{probnum}
\begin{problembox}
\textbf{\theprobnum.} Determine all continuously differentiable functions $f : \mathbb{R} \to \mathbb{R}$ satisfying for all $x \in \mathbb{R}$ the equation
\[
f(x)^2 = \int_0^x f(t)^2 + f'(t)^2\,dt + 2023.
\]
\end{problembox}

\stepcounter{probnum}
\begin{problembox}
\textbf{\theprobnum.} Determine all functions $f : \mathbb{R} \to \mathbb{R}$ satisfying for all $x, y \in \mathbb{R}$ the equation
\[
f(x^2 - y^2) = (x-y)(f(x) + f(y)).
\]
\end{problembox}

\stepcounter{probnum}
\begin{problembox}
\textbf{\theprobnum.} Consider a continuous function $f : \mathbb{R} \to \mathbb{R}$ satisfying
\[
f(2x - f(x)) = x \quad \text{for all } x \in \mathbb{R}.
\]
Suppose that $f(\alpha) = \alpha$ for some $\alpha \in \mathbb{R}$. Show that then $f(x) = x$ for all $x \in \mathbb{R}$. Does this conclusion still hold if $f$ such an $\alpha$ need not necessarily exist?
\end{problembox}

\stepcounter{probnum}
\begin{problembox}
\textbf{\theprobnum.} Let $c > 0$. Determine all continuous functions $f : \mathbb{R} \to \mathbb{R}$ with the property that
\[
f(x) = f(x^2 + c) \quad \text{for all } x \in \mathbb{R}.
\]
\end{problembox}

\stepcounter{probnum}
\begin{problembox}
\textbf{\theprobnum.} Determine all $\alpha > 0$ with the property that for all differentiable functions $f : \mathbb{R} \to \mathbb{R}_{>0}$ such that $f'(x) > f(x)$ for all $x \in \mathbb{R}$, there exists some $N \in \mathbb{R}$ such that $f(x) > e^{\alpha x}$ for all $x > N$.
\end{problembox}

\stepcounter{probnum}
\begin{problembox}
\textbf{\theprobnum.} Let $a, b \in (0,1/2)$ and let $g$ be a continuous function satisfying the property that
\[
g(g(x)) = ag(x) + bx \quad \text{for all } x \in \mathbb{R}.
\]
Prove that $g(x) = cx$ for some $c \in \mathbb{R}$.
\end{problembox}

\stepcounter{probnum}
\begin{problembox}
\textbf{\theprobnum.} Prove that there exists no function $f : \mathbb{R} \to \mathbb{R}$ satisfying
\[
f(f(x) + y) = f(x) + 3x + y f(y) \quad \text{for all } x, y \in \mathbb{R}.
\]
\end{problembox}

\stepcounter{probnum}
\begin{problembox}
\textbf{\theprobnum.} Determine all differentiable functions $f : (0,\infty) \to (0,\infty)$ for which there exists some $a \in \mathbb{R}$ such that
\[
f'(a/x) = x / f(x) \quad \text{for all } x > 0.
\]
\end{problembox}

\stepcounter{probnum}
\begin{problembox}
\textbf{\theprobnum.} Determine all functions $f : \mathbb{R} \setminus \{0\} \to \mathbb{R} \setminus \{0\}$ satisfying
\[
x f(x^2 f(y)) = y f(x) \quad \text{for all } x, y \in \mathbb{R}.
\]
\end{problembox}

\stepcounter{probnum}
\begin{problembox}
\textbf{\theprobnum.} Determine all twice differentiable functions $f : \mathbb{R} \to \mathbb{R}$ with $f'(x) f''(x) = 0$ for all $x \in \mathbb{R}$.
\end{problembox}

\stepcounter{probnum}
\begin{problembox}
\textbf{\theprobnum.} Determine all continuous functions $f : \mathbb{R} \to \mathbb{R}$ satisfying
\[
f(x) + \int_0^x (x-t) f(t)\,dt = 1 \quad \text{for all } x \in \mathbb{R}.
\]
\end{problembox}

\stepcounter{probnum}
\begin{problembox}
\textbf{\theprobnum.} Let $f : [-2,2] \to \mathbb{R}$ be continuously differentiable. Prove that there exists some $x \in (-2,2)$ such that
\[
f'(x) - f(x)^2 < 1.
\]
\end{problembox}

\stepcounter{probnum}
\begin{problembox}
\textbf{\theprobnum.} (*) Let $f : \mathbb{R} \to [0,\infty)$ be a continuously differentiable function. Prove that
\[
\left| \int_0^1 f(x)^3\,dx - f(0)^2 \int_0^1 f(x)\,dx \right| \le \max_{0 \le x \le 1} |f'(x)| \left( \int_0^1 f(x)\,dx \right)^2.
\]
\end{problembox}

\stepcounter{probnum}
\begin{problembox}
\textbf{\theprobnum.} Let $f : (0,\infty) \to \mathbb{R}$ be a twice continuously differentiable function such that
\[
|f''(x) + 2x f'(x) + (x^2+1) f(x)| \le 1
\]
for all $x > 0$. Prove that $\lim_{x \to \infty} f(x) = 0$.
\end{problembox}

\stepcounter{probnum}
\begin{problembox}
\textbf{\theprobnum.} For any given $x \in (0,\pi/2)$ determine which of $\tan(\sin(x))$ and $\sin(\tan(x))$ is bigger.
\end{problembox}

\stepcounter{probnum}
\begin{problembox}
\textbf{\theprobnum.} Let $f : \mathbb{R} \to \mathbb{R}$ be a twice differentiable function satisfying $f(0) = 1$ and $f'(0) = 0$, and further
\[
f''(x) - 5f'(x) + 6f(x) \ge 0 \quad \text{for all } x \ge 0.
\]
Show that for all $x \ge 0$, it holds that
\[
f(x) \ge 3e^{2x} - 2e^{3x}.
\]
\end{problembox}

\stepcounter{probnum}
\begin{problembox}
\textbf{\theprobnum.} Let $0 < a < b$. Prove that
\[
\int_a^b (x^2+1) e^{-x^2}\,dx \ge e^{-a^2} - e^{-b^2}.
\]
\end{problembox}

\stepcounter{probnum}
\begin{problembox}
\textbf{\theprobnum.} Let $f : \mathbb{R} \to \mathbb{R}$ be a continuously differentiable function with the property that
\[
f'(t) > f(f(t)) \quad \text{for all } t \in \mathbb{R}.
\]
Show that
\[
f(f(f(t))) \le 0 \quad \text{for all } t \ge 0.
\]
\end{problembox}

\stepcounter{probnum}
\begin{problembox}
\textbf{\theprobnum.} Let $f : \mathbb{R} \to \mathbb{R}$ be a twice differentiable function with the property that $f(0) = 0$. Show that there exists some $c \in (-\pi/2, \pi/2)$ such that
\[
f''(c) = f(c)(1 + 2\tan(c)^2).
\]
\end{problembox}

\stepcounter{probnum}
\begin{problembox}
\textbf{\theprobnum.} (*) Define $f(x) = \sin(x)/x$. Show that for all $x > 0$ and $n \ge 1$, it holds that
\[
|f^{(n)}(x)| < \frac{1}{n+1},
\]
where $f^{(n)}$ denotes the $n$th derivative of $f$.
\end{problembox}

\stepcounter{probnum}
\begin{problembox}
\textbf{\theprobnum.} Consider the set of continuous functions $f : [0,1] \to \mathbb{R}$ satisfying
\[
f(x) + f(y) \ge |x-y| \quad \text{for all } x, y \in [0,1].
\]
Find the minimal value of
\[
\int_0^1 f(x)\,dx
\]
for functions in this set.
\end{problembox}

\stepcounter{probnum}
\begin{problembox}
\textbf{\theprobnum.} Consider a differentiable function $f : \mathbb{R} \to (0,\infty)$ with the property that for some constant $L > 0$, it holds that
\[
|f'(x) - f'(y)| \le L|x-y| \quad \text{for all } x, y \in \mathbb{R}.
\]
Show that for all $x \in \mathbb{R}$, it holds that
\[
f'(x)^2 < 2Lf(x).
\]
\end{problembox}

\stepcounter{probnum}
\begin{problembox}
\textbf{\theprobnum.} Find all differentiable functions $f : (0,\infty) \to \mathbb{R}$ with the property that
\[
f(b) - f(a) = (b-a) f'(\sqrt{ab}) \quad \text{for all } a, b > 0.
\]
\end{problembox}

\stepcounter{probnum}
\begin{problembox}
\textbf{\theprobnum.} Let $f : (-1,1) \to \mathbb{R}$ be a twice differentiable function such that $2f'(x) + xf''(x) \ge 1$ for all $x \in (-1,1)$. Prove that
\[
\int_{-1}^1 x f(x)\,dx \ge \frac{1}{3}.
\]
\end{problembox}

\stepcounter{probnum}
\begin{problembox}
\textbf{\theprobnum.} Find all twice continuously differentiable functions $f : \mathbb{R} \to (0,\infty)$ such that
\[
f''(x) f(x) \ge 2f'(x)^2 \quad \text{for all } x \in \mathbb{R}.
\]
\end{problembox}

\stepcounter{probnum}
\begin{problembox}
\textbf{\theprobnum.} Find all functions $f : \mathbb{R} \to \mathbb{R}$ that have a continuous second derivative and for which the equality
\[
f(7x+1) = 49 f(x)
\]
holds for all $x \in \mathbb{R}$.
\end{problembox}

\stepcounter{probnum}
\begin{problembox}
\textbf{\theprobnum.}  If $f$ has continuous derivative on $[-1,1]$, find
\[
\lim_{h \to 0^+} \int_{-1}^{1} \frac{h\,f(x)}{h^2+x^2}\,dx,
\]
assuming the limit exists.
\end{problembox}

\stepcounter{probnum}
\begin{problembox}
\textbf{\theprobnum.} Let $f : [0,1] \to \mathbb{R}$ be twice differentiable on $[0,1]$ with $f''(x) > 0$ on $[0,1]$. Show that
\[
4 \int_{1/2}^{1} f(x)\,dx \le f(1) + \int_0^1 f(x)\,dx.
\]
\end{problembox}

\stepcounter{probnum}
\begin{problembox}
\textbf{\theprobnum.} If $f$ has continuous third derivative on $[a,b]$ and has at least three roots in $[a,b]$, show that
\[
\max_{x \in [a,b]} |f(x)| \le \frac{(b-a)^3}{6} \max_{x \in [a,b]} |f^{(3)}(x)|.
\]
\end{problembox}

\end{document}
